% ---------------------------------------------------------------------
%  미리보기 전용 래퍼 (kotex이 없는 환경에서 PDF를 뽑기 위한 파일)
%
%  kotex이 설치된 환경에서는 build-sections-ko.tex 을 쓴다. 이 파일은
%  kotex/xeCJK를 쓰지 않고 fontspec만으로 한글을 조판하므로, acmart의
%  기본 라틴 서체(Libertine)가 CJK 서체로 대체된다. 내용 확인용이며
%  최종 제출본 조판에는 사용하지 않는다.
%
%  컴파일:
%    xelatex build-sections-ko-preview && bibtex build-sections-ko-preview \
%      && xelatex build-sections-ko-preview && xelatex build-sections-ko-preview
% ---------------------------------------------------------------------

\documentclass[sigconf,nonacm]{acmart}

\usepackage{fontspec}
\IfFontExistsTF{Noto Serif CJK KR}{%
  \setmainfont{Noto Serif CJK KR}%
  \setsansfont{Noto Sans CJK KR}%
}{%
  \setmainfont{Noto Serif CJK JP}%
  \setsansfont{Noto Sans CJK JP}%
}
\XeTeXlinebreaklocale "ko"
\XeTeXlinebreakskip = 0pt plus 0.1em

\usepackage{booktabs}
\usepackage{graphicx}
\usepackage{amsmath}
\usepackage{array}
\usepackage{url}

\settopmatter{printacmref=false,printfolios=true}
\setcopyright{none}

\makeatletter
\let\proposition\relax \let\endproposition\relax
\makeatother
\theoremstyle{acmplain}
\newtheorem{proposition}{명제}[section]

% 캡션 이름을 본문 표기(그림/표)와 일치시킨다.
\AtBeginDocument{%
  \renewcommand{\figurename}{그림}%
  \renewcommand{\tablename}{표}%
}


\begin{document}

\title{[초안] 방법과 실험 절 조판 확인용}
\author{저자명}
\affiliation{%
  \institution{소속기관}
  \city{서울}
  \country{대한민국}}
\email{author@example.com}
\renewcommand{\shortauthors}{저자명}

\maketitle

% =====================================================================
%  3장 방법 / 4장 실험 (한국어)
%  acmart [sigconf,nonacm] + kotex, XeLaTeX 또는 LuaLaTeX으로 컴파일
%
%  본문에 삽입: % =====================================================================
%  3장 방법 / 4장 실험 (한국어)
%  acmart [sigconf,nonacm] + kotex, XeLaTeX 또는 LuaLaTeX으로 컴파일
%
%  본문에 삽입: % =====================================================================
%  3장 방법 / 4장 실험 (한국어)
%  acmart [sigconf,nonacm] + kotex, XeLaTeX 또는 LuaLaTeX으로 컴파일
%
%  본문에 삽입: \input{sections-method-experiments-ko}
%  필요 패키지: kotex, booktabs, graphicx, amsmath, array, url
%  그림 경로: ../figures_v2/  (scripts/make_paper_figures.py 산출물)
%  추가 참고문헌: paper-refs-additional.bib 병합 필요
%
%  [선택] 정리 환경 이름을 "Proposition"에서 "명제"로 바꾸려면 전처리부에
%  아래 세 줄을 추가한다. acmart가 이미 proposition을 정의하므로 먼저 해제한다.
%
%    \makeatletter
%    \let\proposition\relax \let\endproposition\relax
%    \makeatother
%    \theoremstyle{acmplain}
%    \newtheorem{proposition}{명제}[section]
%
%  추가하지 않으면 "Proposition 3.1"로 조판되며, 본문의 상호참조는 그대로
%  동작한다.
% =====================================================================

\section{방법}
\label{sec:method}

\subsection{설정: 대칭 추정}
\label{sec:setup}

투자자 유형 집합을 $\mathcal{I}=\{\text{외국인},\text{기관},\text{개인}\}$으로
둔다. 각 유형은 관측된 일별 수급이 해석 가능한 시장 상태 특징 위의 잠재 보상을
드러내는 하나의 에이전트다.

\paragraph{연속 행동}
유형 $i$의 $t$일 총매수금액과 총매도금액을 각각 $V^{buy}_{i,t}$,
$V^{sell}_{i,t}$라 하면 행동을
\begin{equation}
u_{i,t}=\frac{V^{buy}_{i,t}-V^{sell}_{i,t}}
             {V^{buy}_{i,t}+V^{sell}_{i,t}}\in[-1,1]
\label{eq:action}
\end{equation}
로 정의한다. 분모가 시장 전체 거래대금이 아니라 해당 유형의 총거래금액이므로
$u_{i,t}$는 그 유형 내부에서 매수와 매도 중 어느 쪽이 우세했는지를 나타내며,
유형별 거래 규모 차이에 영향받지 않는 공통 척도가 된다. 이진 매수ㆍ매도 라벨
대신 연속 강도를 사용하므로 이진화에서 버려지는 강도 정보가 보존된다.

\paragraph{대칭성}
세 유형은 동일한 상태 특징 $x_{i,t}\in\mathbb{R}^{3}$과 동일한 시장 컨텍스트
$C_t\in\mathbb{R}^{2}$, 그리고 동일한 모형 형태를 부여받는다. 유형별로 다른
특징을 미리 배정하지 않으며, 파라미터 $\theta_i=(\beta_i,B_i,\alpha_i)$만 각자의
데이터로 분리 학습한다. 따라서 결과에 나타나는 유형 간 차이는 설계가 심어둔 것이
아니라 데이터에서 학습된 것이며, 이 대칭성이 유형 간 계수의 직접 비교를
정당화한다.

\paragraph{시점 정렬}
학습 표본은 $(x_{i,t},C_t,a_{i,t+1})$이며 $a_{i,t+1}=u_{i,t+1}$이다. $t$일
종가와 수급으로 구성한 상태가 $t+1$일 행동을 설명하므로 미래 거래는 특징에
사용되지 않는다. 모멘텀과 손실구간 특징은 $t$일 정보를, 군집 특징은 $t-1$일
정보를 사용한다. 예측 대상이 $t+1$일 행동이므로 $t$일 정보도 이용 가능하며,
군집 특징의 추가 지연은 정보 제약이 아니라 동일 거래일 유형 간 수급의 회계적
상쇄를 예측 변수에서 배제하기 위한 설계 선택이다(그림~\ref{fig:crossinv}).

\subsection{보상 특징}
\label{sec:features}

세 특징은 모두 각 검증 분할의 학습 인덱스에서만 적합한 표준화 변환을 거친다.
따라서 이후의 모든 계수는 해당 특징이 학습 구간에서 1표준편차 변할 때의 효과로
읽는다.

\paragraph{(F1) 중기 모멘텀}
조정종가 또는 대체 종가 $P_t$로
\begin{equation}
x^{mom}_{i,t}=\log\!\left(\frac{P_t}{P_{t-20}}\right)
\label{eq:momentum}
\end{equation}
와 같이 계산한다. 20거래일은 약 1개월에 해당하는 사전 지정 창으로, 일별 가격
잡음을 완화하면서도 다음 날 행동과 연결되는 중기 추세를 표현한다. 자료에서
성능을 최적화해 고른 값이 아니라 해석 가능성을 위한 설계값이다. 양의 계수는
추세추종을, 음의 계수는 역추세를 뜻하며, 부호 기대는 유형별 모멘텀 거래
성향에 관한 선행 연구~\cite{grinblatt2000investment,choe1999foreign}에서
도출된다.

\paragraph{(F2) 군집행동}
유형 $i$를 제외한 두 유형 $j,k$의 전일 행동에 대한 횡단면 평균
\begin{equation}
x^{herd}_{i,t}=\frac{u_{j,t-1}+u_{k,t-1}}{2},\qquad j,k\neq i
\label{eq:herd}
\end{equation}
으로 정의한다. 시간평균이 아니라 같은 시점의 다른 두 유형에 대한 평균임에
유의한다. 20일 이동평균을 쓰면 즉각적 반응이 아니라 한 달 누적 수급 압력을
측정하게 되므로, 본 연구는 전일 반응에 초점을 맞춘다. 양의 계수는 다른 유형의
전일 방향을 따르는 행동을, 음의 계수는 반대 방향 행동을 뜻한다. 부호 기대는
기관 군집행동 문헌~\cite{lakonishok1992impact}에서 도출되나, 해당 문헌이 기관
\emph{내부}의 군집을 측정하는 반면 식~\eqref{eq:herd}은 \emph{유형 간} 관계를
측정한다는 점에서 대상이 다르다(\S\ref{sec:construct}).

\paragraph{(F3) 평균단가 대비 손실구간}
유형별 평균단가 대용치 $\bar P_{i,t}$와 당일 총거래 VWAP $P^{vwap}_{i,t}$로
\begin{equation}
x^{uw}_{i,t}=\max\!\left(
\frac{\bar P_{i,t}-P^{vwap}_{i,t}}{\bar P_{i,t}},\,0\right)
\label{eq:underwater}
\end{equation}
와 같이 정의한다. 0에서 절단하므로 이 변수는 거래가격이 평균단가 대용치보다
낮은 구간에서만 작동하며, 그 자체로 준거점 의존적 설계다. 평균단가는 총매수ㆍ
총매도 수량과 금액으로 유효 보유량을 재귀적으로 갱신해 계산하되, 전일 보유량에
망각계수 $\rho=0.98$을 곱해 오래된 거래의 영향을 줄이고 매도는 남은 물량의
평균단가를 바꾸지 않는다. $\rho=0.98$의 반감기는 약 34거래일이고 20거래일 후
약 66.8\%가 남으므로 최근 1--2개월의 거래기억을 강조하며, 이 값 역시 자료에서
최적화하지 않았다. 양의 계수는 평균단가 아래에서 다음 날 매수 강도가 커지는
방향이며, 부호 기대는 처분효과
문헌~\cite{shefrin1985disposition,odean1998investors}에서 도출된다.

이 대용치는 실제 계좌 보유량이 아니라 집계 거래로부터 구성한 근사치다. 같은
유형 안의 서로 다른 계좌가 수행한 매수와 매도가 넷팅되므로 ``대표 투자자의
평균단가''라는 개념 자체가 근사에 불과하며, 따라서 $x^{uw}$의 계수는
처분효과를 직접 검정한 결과로 해석할 수 없다.

\subsection{시장 컨텍스트}
\label{sec:context}

컨텍스트는 두 변수로 구성한다. (C1) KOSPI200 1일 수익률 $r_t$, (C2) 원/달러
환율 수준
\begin{equation}
z_t=\frac{\log e_t-\mu_{t-1,252}}{\sigma_{t-1,252}}
\label{eq:fxcontext}
\end{equation}
이다. 환율의 현재값은 사용하지만 평균과 표준편차는 $t-1$까지의 252거래일만
사용한다. 두 컨텍스트 역시 각 분할의 학습 인덱스에서만 다시 표준화하므로,
계수는 해당 컨텍스트가 학습 구간에서 1표준편차 변할 때의 효과로 읽는다.

컨텍스트와 보상 특징은 서로 독립이 아니다. 삼성전자가 KOSPI200에서 차지하는
비중 때문에 종목 모멘텀과 시장수익률이 겹치며, 특징--컨텍스트 상관은 $-0.29$에서
$0.22$ 범위에 있다(그림~\ref{fig:featcorr}). 따라서 개별 계수는 다른 항이
통제된 조건부 연관이며 단독으로 인과적 크기로 읽어서는 안 된다.

\subsection{근시안적 보상, 정책, 그리고 Lasso로의 환원}
\label{sec:model}

\paragraph{잠재 행동점수}
유형 $i$의 시점별 행동점수를
\begin{equation}
q_{i,t}=(\beta_i+B_iC_t)^{\top}x_{i,t}+\alpha_i^{\top}C_t
\label{eq:score}
\end{equation}
로 둔다. $\beta_i\in\mathbb{R}^{3}$은 컨텍스트가 평균일 때의 기본 보상 가중치,
$B_i\in\mathbb{R}^{3\times 2}$는 특징--컨텍스트 상호작용,
$\alpha_i\in\mathbb{R}^{2}$는 컨텍스트의 직접 주효과다. $\alpha$항은 종목
특징이 평균값일 때도 시장 국면이 행동 수준을 이동시키는 수준효과이고, $B$항은
국면에 따라 각 특징에 대한 민감도가 달라지는 기울기 효과다. 상호작용항을
포함하면서 하위차수 항 $C_t$를 누락하는 것은 표준적 사양 오류이므로,
$\alpha$의 포함은 자료로 정당화할 대상이 아니라 사양상 요구되는 것이다.

\paragraph{근시안적 오목 보상}
연속 행동 $a$에 대한 단기 보상을
\begin{equation}
R_i(x,C,a)=q_{i,t}\,a-\frac{\kappa}{2}a^{2}
\label{eq:reward}
\end{equation}
로 둔다. 이차항은 큰 순매수와 순매도에 대칭적 비용을 부과하고 보상을 행동에
대해 오목하게 만들어 내부해를 가능하게 한다. 이 항이 없으면 $[-1,1]$에서의
최대화 해가 $q$의 부호에 따라 거의 항상 $\pm1$이 되어 행동 \emph{강도}를
재현할 수 없다. 식~\eqref{eq:reward}을 $a$에 대해 최대화하면 닫힌형 정책
\begin{equation}
a^{*}_{i,t+1}=\mathrm{clip}\!\left(\frac{q_{i,t}}{\kappa},-1,1\right)
\label{eq:policy}
\end{equation}
을 얻는다. 곡률 $\kappa$는 식별을 위해 1로 고정한다. $\kappa=2$이면 같은 $q$에서
행동 크기가 $1/2$이 되지만 재학습 시 $q$의 계수를 2배로 바꾸면 같은 행동이
재현되므로, $\kappa=1$은 보상과 계수의 단위를 고정하는 정규화이지 곡률이
경험적으로 1이라는 주장이 아니다.

\paragraph{추정}
유형별 학습목적은 학습집합 $\mathcal{D}_i$에서
\begin{equation}
\min_{\theta_i}\ \frac{1}{|\mathcal{D}_i|}
\sum_{t\in\mathcal{D}_i}\bigl(a^{*}_{i,t+1}-a_{i,t+1}\bigr)^{2}
+\lambda_i\lVert\theta_i\rVert_1
\label{eq:loss}
\end{equation}
이며, Adam~\cite{kingma2015adam}으로 유형별 모형을 독립 학습한다.

\begin{proposition}[Lasso 환원]
\label{prop:lasso}
$z_{i,t}=[\,x_{i,t};\ \mathrm{vec}(x_{i,t}C_t^{\top});\ C_t\,]\in\mathbb{R}^{11}$,
$\theta_i=[\,\beta_i;\ \mathrm{vec}(B_i);\ \alpha_i\,]$로 두면
식~\eqref{eq:score}는 $q_{i,t}=\theta_i^{\top}z_{i,t}$이다. 학습 표본 전체에서
$|q_{i,t}|\le 1$이면 식~\eqref{eq:policy}의 clip이 작동하지 않아 $a^{*}=q$이고,
식~\eqref{eq:loss}은
\begin{equation}
\min_{\theta_i}\ \frac{1}{|\mathcal{D}_i|}
\sum_{t}\bigl(\theta_i^{\top}z_{i,t}-a_{i,t+1}\bigr)^{2}
+\lambda_i\lVert\theta_i\rVert_1
\label{eq:lasso}
\end{equation}
이 된다. 이는 다음 날 순매수 강도를 확장 상태특징 $z$에 회귀하는 $L_1$
회귀(Lasso)와 정확히 같다.
\end{proposition}

즉 $\kappa=1$의 근시안적 오목 보상 아래에서 ``선호를 복원한다''와 ``수급을
회귀한다''는 수치적으로 구별되지 않는다. 그림~\ref{fig:saturation}은
명제~\ref{prop:lasso}의 전제가 본 표본에서 실제로 성립함을 보인다. 세 유형 모두
잠재 행동점수의 분포가 행동 경계 $\pm1$로부터 크게 떨어져 있으며($|q|$의
최대값은 외국인 0.842, 개인 0.587, 기관 0.135), 45개 분할 전부에서 포화 비율이
정확히 0이다.

\paragraph{이 환원을 숨기지 않는 이유}
이 일치는 방법의 약점이 아니라 근시안 가정의 논리적 귀결이다. 에이전트가 한
걸음만 내다본다고 가정하면 오늘의 행동이 내일 이후에 영향을 주지 않으므로 내다볼
미래가 없고, 따라서 무엇을 원하는가(보상)와 무엇을 하는가(정책)가 하나로
겹친다. 최대엔트로피 IRL~\cite{ziebart2008maximum}의 확률적 정책
$\pi(a\mid x,C)\propto\exp(R/\tau)$를 식~\eqref{eq:reward}에 적용하면 $\pi$는
평균 $q$, 분산 $\tau$의 절단정규분포가 되어 로그가능도 최대화의 점추정치가
식~\eqref{eq:lasso}과 일치하므로, 이 환원은 특정 구현의 우연이 아니라
근시안ㆍ이차보상 설정 전반의 성질이다.

역강화학습 정식화의 구별되는 가치는 확장 경로에 있다. 상태 동역학과 다기간
가치를 도입하는 순간 보상은 정책과 갈라지고, 회귀로는 도달할 수 없는 구조적
복원으로 이어진다. 본 논문은 그 경로의 첫 단계이며, 성립 조건을 명시한 위에
해석과 검증을 더한다.

\subsection{검증 프로토콜}
\label{sec:protocol}

복원된 보상을 계수 부호와 크기의 보고로 끝내지 않고 다음 네 단계를 적용한다.

\begin{description}
\item[(V1) CPCV 표본외 안정성] 10개 시간 폴드 중 2개를 시험 폴드로 선택하는
조합형 정화 교차검증~\cite{lopez2018advances}을 사용한다. 조합 수는
$\binom{10}{2}=45$이며, 학습과 시험 사이에 1일 purge와 5일 embargo를 적용하고
모든 표준화 변환은 각 분할의 학습 인덱스에만 적합한다. 45개 분할에서 지배 부호가
나타난 분할의 비율(부호 일관성)을 보고하되, 분할들이 관측을 공유하므로 부호
일관성은 독립표본 $p$-값이 아니라 안정성 지표로만 해석한다.

\item[(V2) 달력월 블록 부트스트랩] 부호 일관성이 검정이 아니라는 한계를
보완하기 위해, 표본을 달력월 단위 블록으로 나누고 월 블록을 복원추출해 1{,}000회
재추정하여 계수의 95\% 신뢰구간을 구한다. 월 블록을 쓰는 이유는 일별 수급의
자기상관과 월중 패턴을 블록 내부에 보존하기 위해서다.

\item[(V3) Expanding walk-forward] CPCV는 시험 폴드가 학습 폴드보다 앞설 수
있으므로, 시간 방향을 엄격히 지키는 검증을 추가한다. 최초 400개 관측으로 학습해
5일 purge 후 다음 60거래일을 예측하고, 학습 구간을 60거래일씩 확장하며
반복한다(총 10개 창).

\item[(V4) 정규화 방식 교체] 결과가 $L_1$ 벌점의 산물이 아님을 확인하기 위해
동일 프로토콜을 $L_2$ 벌점과 무벌점으로 재실행하고 계수 부호를 비교한다.
\end{description}

\paragraph{평가지표}
행동 강도 오차는 MAE와 RMSE로, 방향 재현은
$\mathrm{sign}(a^{*})=\mathrm{sign}(a)$의 비율인 방향 정확도로 평가한다. 실제
행동과 예측 행동의 피어슨 상관, $|a^{*}|=1$인 포화 비율을 함께 보고한다. 유형 간
행동 표준편차가 다르므로 MAE와 RMSE의 절대값은 유형 간 비교에 쓰지 않고 각 유형
안에서 해석하며, 표준편차로 정규화한 오차를 병기한다. 참조 기준선으로 전일
자기행동을 그대로 예측하는 지속성 규칙을 함께 보고한다(\S\ref{sec:notclaim}).

\paragraph{구성 타당성 점검}
복원된 부호를 독립적으로 확립된 행동재무 규칙성에서 도출한 사전 기대와
대조한다(\S\ref{sec:construct}). 이 대조는 사후 해석이 아니라 사전에 정한 부호
기대와의 비교다.

% =====================================================================
\section{실험}
\label{sec:experiments}

\subsection{데이터와 설정}
\label{sec:data}

원자료는 삼성전자(005930) 일별 자료이며 가격, 거래대금, 유형별 총매수ㆍ총매도
금액과 수량, 원/달러 환율, KOSPI200 수익률을 포함한다. 252일 환율 이력, 20일
모멘텀, 다음 날 행동을 모두 계산할 수 있는 관측 중 최신 973행을 분석 표본으로
사용하였다. 상태일은 2022년 1월 6일--2025년 12월 29일이며, 마지막 상태는
2025년 12월 30일 행동을 목표로 한다.

단일 종목은 표본의 한계가 아니라 교차종목 교란 없이 유형 간 선호 이질성만
분리하기 위한 통제된 시연 설정이다. 종목 고정효과, 유동성 차이, 지수 편입 효과가
개입하지 않으므로 유형 간 계수 차이를 설계가 아닌 학습의 결과로 귀속할 수 있다.

\begin{table}[t]
  \caption{데이터 및 추정 설정}
  \label{tab:setup}
  \small
  \begin{tabular}{@{}>{\raggedright\arraybackslash}p{0.30\columnwidth}
                  >{\raggedright\arraybackslash}p{0.62\columnwidth}@{}}
    \toprule
    항목 & 설정 \\
    \midrule
    표본 & 973거래일, 2022-01-06--2025-12-29 \\
    유형 & 외국인, 기관, 개인 (대칭 추정) \\
    보상 특징 & 모멘텀(20일), 군집(전일), 손실구간($\rho=0.98$) \\
    컨텍스트 & KOSPI200 1일 수익률, USD/KRW 252일 수준 $z$ \\
    유형별 파라미터 & 11개 ($\beta$ 3, $B$ 6, $\alpha$ 2) \\
    공통 설정 & seed 42, CPU, $\kappa=1$, 학습구간 전용 표준화 \\
    외국인 & 75 epochs, batch 256, lr 0.001, $\lambda=0$ \\
    기관 & 10 epochs, batch 2048, lr 0.0005, $\lambda=0.01$ \\
    개인 & 20 epochs, batch 512, lr 0.001, $\lambda=0.0003$ \\
    검증 & CPCV 10 folds, test 2, purge 1, embargo 5, 45 splits \\
    안정성 & 월 블록 부트스트랩 1{,}000회, walk-forward 10창, 정규화 교체 \\
    \bottomrule
  \end{tabular}
\end{table}

행동 평균은 외국인 $-0.0063$, 기관 $-0.0100$, 개인 $0.0035$로 세 유형 모두 0에
가까우나 표준편차는 개인(0.334)이 가장 크고 기관(0.125)이 가장 작다. 1차
자기상관은 외국인 0.401, 기관 0.149, 개인 0.345다.

동일 거래일 유형 간 행동은 강한 음의 관계를 갖는다(그림~\ref{fig:crossinv}).
외국인--개인 상관이 $-0.772$, 기관--개인이 $-0.499$이며, 이는 한 유형의 매수가
다른 유형의 매도로 나타날 수밖에 없는 시장 청산 제약을 반영한다. 이 사실은
\S\ref{sec:beta}과 \S\ref{sec:alpha}의 해석 범위를 제한하는 근거로 다시
사용된다.

\subsection{표본외 행동 재현}
\label{sec:oos}

\begin{figure}[t]
  \centering
  \includegraphics[width=\columnwidth]{../figures_v2/fig2_actual_vs_prediction.pdf}
  \caption{투자자별 실제 행동과 CPCV 표본외 예측 행동. 각 날짜의 9개 CPCV
  시험예측을 먼저 평균한 뒤 실제값과 예측값에 20거래일 이동평균을 적용했다.
  괄호 안 상관은 이동평균 계열 간 값이며 표~\ref{tab:performance}의 일별
  표본외 상관과 직접 비교하지 않는다.}
  \Description{외국인, 기관, 개인의 실제 순매수 행동과 모형 예측을 시간축 위에
  겹쳐 그린 세 개의 선그래프. 기관 패널에서 예측선이 0 근방에 평탄하게 놓여 있다.}
  \label{fig:prediction}
\end{figure}

그림~\ref{fig:prediction}은 세 유형의 실제 순매수 행동과 표본외 예측 행동을
나란히 보인다. 세 가지가 한눈에 읽힌다. 첫째, 외국인과 개인에서는 예측 계열이
실제 계열의 국면 전환을 따라간다(이동평균 계열 간 상관 0.69, 0.76). 둘째, 예측
진폭이 실제보다 뚜렷하게 작으며, 이는 $L_1$ 정규화에 의한 계수 축소의 결과다.
셋째, 기관의 예측 계열은 사실상 0선에 붙어 있다(0.02). \S\ref{sec:null}에서
보고하는 null이 계수 표가 아니라 예측 시계열에서 먼저 드러난다.

\begin{table}[t]
  \caption{CPCV 표본외 성능(45개 분할 평균 $\pm$ 표준편차)과 지속성 기준선}
  \label{tab:performance}
  \small
  \centering
  \begin{tabular}{@{}lccc@{}}
    \toprule
    지표 & 외국인 & 기관 & 개인 \\
    \midrule
    \multicolumn{4}{@{}l}{\emph{주모형}}\\
    방향 정확도   & $0.636\pm0.039$ & $0.542\pm0.038$ & $0.608\pm0.044$ \\
    상관계수      & $0.278\pm0.110$ & $0.068\pm0.064$ & $0.245\pm0.118$ \\
    MAE           & $0.195\pm0.019$ & $0.094\pm0.008$ & $0.264\pm0.029$ \\
    RMSE          & $0.246\pm0.021$ & $0.125\pm0.012$ & $0.318\pm0.030$ \\
    MAE / 행동 SD & $0.759$ & $0.756$ & $0.791$ \\
    포화 비율     & $0.000$ & $0.000$ & $0.000$ \\
    \addlinespace
    \multicolumn{4}{@{}l}{\emph{지속성 기준선} ($a_{t+1}\leftarrow a_t$)}\\
    방향 정확도   & $\mathbf{0.642}$ & $\mathbf{0.547}$ & $\mathbf{0.616}$ \\
    상관계수      & $\mathbf{0.359}$ & $\mathbf{0.136}$ & $\mathbf{0.311}$ \\
    MAE           & $0.223$ & $0.123$ & $0.299$ \\
    \addlinespace
    행동의 1차 자기상관 & $0.401$ & $0.149$ & $0.345$ \\
    \bottomrule
  \end{tabular}
\end{table}

\begin{figure}[t]
  \centering
  \includegraphics[width=\columnwidth]{../figures_v2/fig5_oos_performance.pdf}
  \caption{표본외 성능의 분할 간 분포와 자명한 기준선. 실선은 주모형 평균,
  파선은 지속성 기준선이다. 세 유형 모두에서 파선이 실선보다 위에 있다.}
  \Description{방향 정확도와 상관계수에 대해 세 투자자 유형의 45개 분할 값을
  바이올린 및 점으로 그린 두 개의 패널.}
  \label{fig:oos}
\end{figure}

표~\ref{tab:performance}과 그림~\ref{fig:oos}에서 평균만 보면 가려지는 두 가지가
드러난다. 첫째, 분할 간 산포가 크다. 특히 상관계수는 외국인이 $-0.05$에서
$0.45$까지 걸쳐 있어 단일 분할의 수치로 유형을 비교하는 것은 의미가 없다. 둘째,
자명한 지속성 기준선이 방향 정확도와 상관계수에서 주모형을 앞선다. $a_t$는 $t$일
종가 이후 공개되므로 이는 정보상 이용 가능한 규칙이다. 주모형이 앞서는 것은
MAE뿐인데, 이는 정규화로 예측 진폭이 축소되어 크기 오차가 줄어든 결과이므로
정보량 우위로 읽을 수 없다.

이 격차의 원인은 명확하다. 행동의 1차 자기상관이 $0.401/0.149/0.345$인데 본
연구의 보상 특징 집합은 자기 지연행동을 포함하지 않는다. 특징을 (F1)--(F3)로
제한한 것은 세 유형이 독립적으로 확립된 행동재무 규칙성에 대응하는 특징만
공유하도록 설계했기 때문이며, 예측 성능을 극대화하려는 설계가 아니다. 본 연구는
이 결과를 은폐하지 않고 표~\ref{tab:performance}에 함께 싣는다.
표~\ref{tab:performance}과 그림~\ref{fig:oos}는 모형이 잡음만 학습한 것이 아님을
보이는 용도로 사용되며, 예측 성능의 우위를 주장하는 근거로 사용되지 않는다.

모든 유형ㆍ모든 분할에서 포화 비율이 정확히 0이므로 명제~\ref{prop:lasso}의 전제
$|q|\le1$이 표본 전체에서 성립한다. 즉 본 추정은 실제로 식~\eqref{eq:lasso}의
Lasso와 동일하며, 성능이 $\pm1$ 극단 예측에 의존하지도 않는다.

\subsection{복원된 보상: 모멘텀 축의 거울상}
\label{sec:beta}

\begin{figure*}[t]
  \centering
  \includegraphics[width=\textwidth]{../figures_v2/fig1_reward_weights_bar.pdf}
  \caption{투자자별로 복원된 기본 보상가중치 $\beta$. 파란색은 양, 붉은색은 음의
  평균 계수이며, 오차막대는 45개 분할의 1 표준편차, 괄호 안은 부호 일관성이다.}
  \Description{외국인, 기관, 개인 각각에 대해 모멘텀, 군집, 손실구간 계수를 그린
  세 개의 막대그래프. 외국인 모멘텀은 양, 개인 모멘텀은 음이며 기관의 세 막대는
  거의 보이지 않는다.}
  \label{fig:betabar}
\end{figure*}

\begin{figure*}[t]
  \centering
  \includegraphics[width=\textwidth]{../figures_v2/fig3_beta_alpha_forest.pdf}
  \caption{보상가중치와 컨텍스트 주효과의 분할 간 안정성. 회색 점은 개별 CPCV
  분할, 굵은 오차선은 45개 분할의 표준편차, 옅은 가로선은 달력월 블록 부트스트랩
  95\% 신뢰구간이다. 속이 빈 표식은 신뢰구간이 0을 포함하는 항이다.}
  \Description{패널 A는 기본 보상가중치, 패널 B는 컨텍스트 주효과에 대한 forest
  plot으로, 각 항의 분할별 추정치와 신뢰구간을 보인다.}
  \label{fig:forest}
\end{figure*}

그림~\ref{fig:betabar}는 세 유형에 대해 복원된 기본 보상가중치를 나란히 보인다.
동일한 특징 집합과 동일한 모형 형태를 부여하고 파라미터만 분리 학습했으므로, 세
패널의 높이 차이는 전적으로 데이터에서 학습된 것이다.
그림~\ref{fig:forest}는 같은 $\beta$와 컨텍스트 주효과 $\alpha$를 분할 수준까지
펼쳐 안정성을 보인다.

\paragraph{모멘텀 축의 거울상}
외국인의 모멘텀 가중치는 $+0.0496$으로 45개 분할 전부에서 양수였고, 개인은
$-0.0304$로 45개 분할 전부에서 음수였다. 그림~\ref{fig:betabar}에서 두 막대가
0선을 사이에 두고 반대 방향을 향하며, 그림~\ref{fig:forest}에서 두 부트스트랩
신뢰구간은 서로 겹치지 않고 각각 0을 제외한다. 동일한 특징 집합을 유형별로 독립
추정했는데도 두 유형이 모멘텀에서 정반대로 갈린다. 외국인은 추세추종, 개인은
역추세로 나타나며, 이는 선행 연구~\cite{grinblatt2000investment,choe1999foreign}의
부호 기대와 일치한다.

\paragraph{군집}
군집 계수는 세 유형 모두 45개 분할 전부에서 음수였고, 외국인($-0.0402$)과
개인($-0.0322$)은 부트스트랩 신뢰구간도 0을 제외한다. 그러나 이 결과는 행동적
해석과 시장 청산 제약을 분리하지 못한다. 그림~\ref{fig:crossinv}가 보이듯 동일
거래일 유형 간 행동 상관이 $-0.77$에 이르므로, 음의 군집 계수 중 상당 부분이
회계적 제약에서 발생할 수 있다. 식~\eqref{eq:herd}가 1거래일 지연값을 쓰므로
동시점 제약이 직접 적용되지는 않으나, 행동의 자기상관을 통해 제약이 전이될 수
있다. 따라서 본 연구는 이를 ``관측된 집계 수급이 유형 간 음의 교차 연관을
보인다''는 서술적 결과로 제한하고, 각 유형이 독립적으로 반대매매 성향을 갖는다는
행동적 주장으로 확장하지 않는다.

\paragraph{손실구간}
개인의 손실구간 가중치는 $+0.0293$으로 45개 분할 전부에서 양수였다. 평균단가
대용치 아래에서 다음 날 매수 강도가 커지는 방향이며, 처분효과
문헌~\cite{shefrin1985disposition,odean1998investors}의 기준가격 의존성과
정합한다. 외국인도 $+0.0094$(95.6\%)로 같은 부호이나 크기가 개인의 $1/3$
수준이다.

다만 그림~\ref{fig:forest}에서 보듯 손실구간은 세 유형 모두 부트스트랩 신뢰구간이
0을 포함한다. 개인의 구간은 $[-0.0122,+0.0665]$로, 45개 분할에서 부호가 100\%
일관됨에도 불구하고 월 블록 재표집에서는 부호가 뒤집히는 경우가 5.5\% 발생한다.
이는 CPCV 분할들이 관측을 공유하기 때문에 부호 일관성이 불확실성을 과소평가한다는
\S\ref{sec:protocol}의 경고가 실제로 작동한 사례다. 따라서 본 연구는 처분효과와
정합하는 부호가 관측되었다고 보고하되, 이를 통계적으로 확립된 결과로 주장하지
않는다.

\subsection{컨텍스트 반응: 거울상의 재현}
\label{sec:alpha}

모멘텀 축에서 관찰된 거울상이 컨텍스트 주효과에서도 그대로
재현된다(그림~\ref{fig:forest}의 패널 B). 외국인은 시장 상승일에 다음 날 순매수
방향이 강해지고($+0.0361$) 원화 약세 국면에서는 순매도 방향으로
이동하는($-0.0261$) 반면, 개인은 두 컨텍스트 모두에서 정확히
반대였다($-0.0246$, $+0.0298$). 네 계수 모두 45개 분할 전부에서 부호가
유지되었고 부트스트랩 신뢰구간도 모두 0을 제외한다. 기관의 두 계수는 부호
일관성이 95\% 이상이나 절대 크기가 외국인ㆍ개인의 $1/8$ 수준이고 신뢰구간이 0을
포함한다.

이 거울상 역시 군집 계수와 같은 해석상 유보가 필요하다.
그림~\ref{fig:crossinv}의 $-0.772$라는 상관 하에서 두 유형의 컨텍스트 반응이
반대 부호로 나타나는 것은 부분적으로 시장 청산 제약을 반영할 수 있다. 본 논문은
이를 두 유형이 시장ㆍ환율 국면에서 반대편 수급을 형성한다는 서술적 사실로
제시한다.

\paragraph{상호작용}
12개 상호작용 계수 중 45개 분할 전부에서 부호가 유지된 것은 외국인의
모멘텀 $\times$ KOSPI200($-0.0139$), 개인의 모멘텀 $\times$ 환율($+0.0144$),
개인의 군집 $\times$ KOSPI200($-0.0080$), 개인의 손실구간 $\times$
환율($-0.0186$) 네 개이며, 이 중 부트스트랩을 통과하는 것은 앞의 하나와 마지막 하나다. 12개 계수를
동시에 검토했으므로 본 연구는 상호작용을 탐색적 결과로만 제시한다.

\subsection{안정성 프로토콜의 결과}
\label{sec:stability}

\paragraph{(V2) 달력월 블록 부트스트랩}
그림~\ref{fig:forest}의 옅은 가로선이 그 결과다. 부트스트랩은 CPCV 부호
일관성보다 엄격한 기준이며, 결과는 셋으로 갈린다. 첫째, 모멘텀 거울상과 컨텍스트
거울상은 부트스트랩에서도 유지된다. 외국인ㆍ개인의 모멘텀, 군집, 두 컨텍스트
주효과 등 8개 항이 모두 0을 제외한다. 둘째, 손실구간 가중치는 통과하지 못한다.
셋째, 기관은 어떤 항도 0을 제외하지 못하며, 그림~\ref{fig:forest}에서 기관의
여섯 표식이 모두 속이 비어 있다.

\begin{figure*}[t]
  \centering
  \includegraphics[width=\textwidth]{../figures_v2/fig6_walk_forward.pdf}
  \caption{Expanding walk-forward 10개 창의 계수 경로와 방향 정확도. 성능은
  창마다 흔들리지만 계수 부호는 유지된다.}
  \Description{확장 창 번호를 가로축으로 하여 세 유형의 모멘텀 계수, 두 컨텍스트
  주효과, 방향 정확도를 그린 네 개의 선그래프.}
  \label{fig:walkforward}
\end{figure*}

\paragraph{(V3) Expanding walk-forward}
그림~\ref{fig:walkforward}가 그 결과다. 두 가지가 동시에 읽힌다. 성능은 창마다
크게 흔들린다. 외국인 방향 정확도는 0.43에서 0.85 사이를 오가며 평균
$0.610\pm0.116$, 상관계수는 $0.131\pm0.162$로 CPCV의 0.278에서 절반 이하로
떨어진다. 개인도 $0.245\rightarrow0.129$다. CPCV는 시험 폴드를 학습 폴드보다
앞세울 수 있는 반면 walk-forward는 그렇지 않으므로, 본 연구는 이 저하를 시간
구조에서 오는 일반화 한계로 그대로 보고한다.

반면 계수 경로는 흔들리지 않는다. 외국인의 모멘텀 계수는 10개 창 전부에서 양수,
개인은 전부에서 음수이며 두 경로가 한 번도 교차하지 않는다. 두 컨텍스트
주효과에서도 같은 거울상이 창 전체에 걸쳐 유지된다. 기관의 세 경로는 모두 0선에
붙어 있다. 즉 성능은 시간에 따라 흔들리지만 복원된 선호의 방향은 흔들리지
않는다.

\begin{figure*}[t]
  \centering
  \includegraphics[width=\textwidth]{../figures_v2/fig7_institution_null.pdf}
  \caption{기관의 수렴된 null. 좌: 복원된 선호의 크기(로그 축). 우: 정규화 방식
  교체에 대한 민감도이며 세 패널은 각각 독립 $y$축을 쓴다.}
  \Description{왼쪽은 유형별 계수 절대값을 로그 축 막대로, 오른쪽은 L1, L2,
  무벌점 하의 계수를 유형별로 그린 세 개의 선그래프.}
  \label{fig:null}
\end{figure*}

\paragraph{(V4) 정규화 방식 교체}
그림~\ref{fig:null}이 그 결과다. 외국인과 개인의 계수는 벌점 방식에 사실상
불변이다(소수점 넷째 자리까지 일치). 따라서 모멘텀 거울상은 $L_1$ 벌점의 산물이
아니다. 반면 기관의 계수만 벌점에 따라 이동하며(군집
$-0.0076\rightarrow-0.0124$, 손실구간 $-0.0000\rightarrow-0.0018$), 그 $y$축
범위는 다른 두 유형의 $1/20$ 수준이다. 계수가 벌점 선택에 좌우된다는 것은
데이터가 그 값을 결정하지 못한다는 뜻이다.

\subsection{기관: 수렴된 null}
\label{sec:null}

기관 수급에 대해 본 연구는 일별 빈도에서 안정적으로 식별되는 선호가 없다는
결과를 얻었고, 이를 은폐하지 않고 발견으로 보고한다. 근거는 다섯이며 모두 그림에
나타난다.

\begin{enumerate}
\item \textbf{예측의 퇴화.} 그림~\ref{fig:prediction} 중간 패널에서 기관의 예측
계열은 실제 계열이 $\pm0.1$ 사이를 오가는 동안 0선에 붙어 있다.
\item \textbf{크기.} 그림~\ref{fig:betabar}에서 기관 패널의 세 막대는 다른 두
유형과 같은 $y$축 위에서 거의 보이지 않는다. 군집을 제외하면 모멘텀 $-0.0006$,
손실구간 $-0.0008$로 $1/40$ 수준이다.
\item \textbf{불확실성.} 그림~\ref{fig:forest}에서 기관의 여섯 표식은 모두 속이
비어 있다.
\item \textbf{벌점 의존성.} 그림~\ref{fig:null} 우측 중앙 패널에서 기관만 벌점에
따라 계수가 이동한다.
\item \textbf{시간 불안정.} 그림~\ref{fig:walkforward}에서 기관의 경로는 0선
근방에 붙어 있으며, 모멘텀 부호 일관성은 10개 창 중 40\%, 손실구간은 0\%다.
\end{enumerate}

\begin{figure*}[t]
  \centering
  \includegraphics[width=\textwidth]{../figures_v2/fig4_training_convergence.pdf}
  \caption{45개 CPCV 분할에서의 학습 수렴. 옅은 회색선은 개별 분할, 굵은 선과
  띠는 평균 $\pm$ 1 표준편차다. 주석은 도달한 해가 동일 목적함수의 폐쇄형
  최적해 대비 남긴 상대 격차다.}
  \Description{에포크를 가로축으로 하여 세 유형의 훈련 MSE 곡선을 그린 세 개의
  패널. 기관 패널의 감소폭이 가장 작다.}
  \label{fig:convergence}
\end{figure*}

이 null은 최적화 실패가 아니다. 그림~\ref{fig:convergence}는 45개 분할의 학습
곡선을 유형별로 보인다. 기관의 훈련 MSE는 학습 구간 전체에서 1.5\%만 감소하며
곡선이 시각적으로 평탄해 보이지 않는다. 그러나 이는 학습 부족이 아니라 애초에
줄일 수 있는 오차가 거의 없기 때문이다. 명제~\ref{prop:lasso}에 따라
식~\eqref{eq:loss}은 볼록 문제이므로 전역 최적해가 유일하며, 도달한 해와 동일
목적함수의 폐쇄형 최적해를 비교하면 기관의 목적함수 상대 격차는 평균 0.17\%,
최대 0.48\%에 불과하다(외국인 0.21\%, 개인 0.92\%). 즉 세 유형 모두 도달 가능한
최적해에 근접했으며, 기관에서만 그 최적해의 계수가 0 근방이다. 데이터가 기관에
대해 결정하는 선호가 없는 것이지, 최적화가 그것을 찾지 못한 것이 아니다.

경제적으로 이 null은 자연스럽게 읽힌다. ``기관''은 연기금, 투신, 보험, 사모,
금융투자를 합산한 범주이므로 서로 다른 목적함수를 가진 주체들의 수급이 상쇄된다.
유일하게 일관된 음의 군집 계수는 이들이 공통적으로 다른 유형의 반대편에 서는
유동성 공급자 역할을 한다는 해석과 정합하나, 부트스트랩 신뢰구간이 경계에
걸치므로 확정하지 않는다.

\subsection{구성 타당성 점검}
\label{sec:construct}

\begin{table}[t]
  \caption{사전 부호 기대와 복원 결과}
  \label{tab:construct}
  \small
  \begin{tabular}{@{}>{\raggedright\arraybackslash}p{0.26\columnwidth}
                  >{\raggedright\arraybackslash}p{0.15\columnwidth}
                  >{\raggedright\arraybackslash}p{0.22\columnwidth}
                  >{\raggedright\arraybackslash}p{0.24\columnwidth}@{}}
    \toprule
    규칙성 & 사전 기대 & 복원 결과 & 판정 \\
    \midrule
    외국인의 추세거래~\cite{grinblatt2000investment,choe1999foreign}
      & $\beta^{mom}>0$ & $+0.0496$ (100\%, CI 0 제외) & 부합 \\
    개인의 역추세 거래~\cite{grinblatt2000investment,kaniel2008individual}
      & $\beta^{mom}<0$ & $-0.0304$ (100\%, CI 0 제외) & 부합 \\
    처분효과~\cite{shefrin1985disposition,odean1998investors}
      & $\beta^{uw}>0$ & $+0.0293$ (100\%, CI 0 포함) & 부분 부합 \\
    기관의 군집행동~\cite{lakonishok1992impact}
      & $\beta^{herd}>0$ & $-0.0044$ (100\%, CI 0 포함) & 불일치 \\
    \bottomrule
  \end{tabular}
\end{table}

표~\ref{tab:construct}이 대조 결과다. 세 유형에 동일한 특징과 동일한 모형 형태를
부여하고 파라미터만 분리 학습했음에도, 복원된 부호가 각각 독립적으로 확립된
규칙성과 부합하고 두 유형이 모멘텀에서 정반대로 갈린다. 이는 복원된 보상이 임의의
적합 결과가 아니라는 구성 타당성의 증거다.

불일치 항목은 측정 대상이 다르기 때문이다. Lakonishok
외~\cite{lakonishok1992impact}는 기관 \emph{내부} 주체들이 같은 방향으로 몰리는
정도를 측정하는 반면, 식~\eqref{eq:herd}은 기관이 \emph{다른 두 유형}의 전일
수급과 같은 방향으로 움직이는지를 측정한다. 따라서 본 결과는 해당 문헌의 반증이
아니며, 유형 내부 군집을 검정하려면 기관 세부 주체별 자료가 필요하다.

\subsection{본 연구가 주장하지 않는 것}
\label{sec:notclaim}

프로토콜의 마지막 단계로, 본 결과가 지지하지 않는 주장을 명시한다.

\begin{enumerate}
\item \textbf{수익률 예측가능성.} 본 연구의 지표는 수급의 방향과 강도에 대한
재현 오차이며 수익률, 샤프지수, 거래비용을 포함하지 않는다.
표~\ref{tab:performance}의 수치로부터 어떤 매매 규칙도 도출되지 않는다.

\item \textbf{수급 예측 성능의 우위.} 표~\ref{tab:performance}과
그림~\ref{fig:oos}에서 보고했듯 지속성 기준선이 방향 정확도와 상관계수에서
주모형을 앞선다. 본 연구의 목적은 예측 성능이 아니라 해석 가능한 공통 척도에서의
선호 복원이며, 두 목적은 다른 특징 집합을 요구한다. 예비 점검에서 지연
자기행동을 네 번째 특징으로 추가하면 모형이 기준선을 넘어서는 반면(방향 정확도
$0.646/0.569/0.629$), 개인의 $\alpha^{KOSPI}$ 부호가 반전되고 손실구간 계수가
소멸한다. 즉 본 논문이 보고한 거울상 중 일부는 수급 자기상관과 분리되지 않는다.

\item \textbf{인과효과.} 복원된 가중치는 표준화된 관측 특징 사이의 조건부
연관이다. 그림~\ref{fig:featcorr}이 보이듯 컨텍스트와 보상 특징 사이에 $-0.29$에서
$0.22$ 범위의 상관이 존재하므로 개별 계수의 크기는 다른 항의 통제 조건에
의존한다.

\item \textbf{유형 간 음의 연관의 행동적 해석.} 군집 계수와 거울상 컨텍스트
반응은 시장 청산 제약과 분리되지 않는다(그림~\ref{fig:crossinv}).

\item \textbf{처분효과의 검정.} 개인의 손실구간 부호는 문헌과 정합하나 월 블록
부트스트랩을 통과하지 못했고(그림~\ref{fig:forest}), 평균단가 대용치는 계좌별
실현손익과 보유기간이 아니라 집계 거래에서 구성한 근사치다.

\item \textbf{종목 간 일반화.} 단일 종목 설정은 유형 간 이질성 분리를 위한
통제이며, 다른 종목과 시장으로의 확장은 검증되지 않았다.
\end{enumerate}

% =====================================================================
\appendix

\section{보충 자료}
\label{sec:appendix}

본문에서 참조한 보조 그림과 표를 모은다. 그림~\ref{fig:saturation}은
명제~\ref{prop:lasso}의 전제를, 그림~\ref{fig:featcorr}은 특징과 컨텍스트의
상관을, 그림~\ref{fig:crossinv}은 유형 간 동시점 행동을 보인다.
표~\ref{tab:interactions}은 본문에서 탐색적 결과로만 제시한 상호작용 계수
전체이고, 표~\ref{tab:convergence}은 도달한 해의 최적성 점검이다.

\begin{figure}[htbp]
  \centering
  \includegraphics[width=\columnwidth]{../figures_v2/figA1_no_saturation.pdf}
  \caption{잠재 행동점수의 분포와 행동 경계. $|q|\le1$이므로
  명제~\ref{prop:lasso}의 전제가 표본 전체에서 성립하며, 45개 분할의 포화 비율은
  정확히 0이다.}
  \Description{세 유형의 잠재 행동점수 히스토그램과 $\pm1$ 경계선.}
  \label{fig:saturation}
\end{figure}

\begin{figure}[htbp]
  \centering
  \includegraphics[width=\columnwidth]{../figures_v2/figA2_feature_correlation.pdf}
  \caption{보상 특징과 컨텍스트의 상관행렬(외국인 특징 기준, 973개 관측).}
  \Description{다섯 변수 사이의 상관계수를 색과 숫자로 표시한 히트맵.}
  \label{fig:featcorr}
\end{figure}

\begin{figure}[htbp]
  \centering
  \includegraphics[width=\columnwidth]{../figures_v2/figA3_cross_investor_actions.pdf}
  \caption{동일 거래일 유형 간 행동. 점선은 완전 상쇄선이며, 외국인과 개인의
  행동이 그 근방에 밀집한다.}
  \Description{유형 쌍별 행동 산점도 세 개와 각 쌍의 상관계수.}
  \label{fig:crossinv}
\end{figure}

\begin{table}[htbp]
  \caption{상호작용 계수 $B$(45개 분할)}
  \label{tab:interactions}
  \small
  \centering
  \begin{tabular}{@{}llrrr@{}}
    \toprule
    유형 & 특징 $\times$ 컨텍스트 & 평균 & 표준편차 & 부호 일관성 \\
    \midrule
    외국인 & 모멘텀 $\times$ KOSPI200 & $-0.0139$ & $0.0052$ & 100.0\% \\
    외국인 & 모멘텀 $\times$ 환율     & $-0.0030$ & $0.0057$ & 71.1\% \\
    외국인 & 군집 $\times$ KOSPI200   & $+0.0016$ & $0.0074$ & 66.7\% \\
    외국인 & 군집 $\times$ 환율       & $+0.0037$ & $0.0073$ & 68.9\% \\
    외국인 & 손실구간 $\times$ KOSPI200 & $-0.0062$ & $0.0040$ & 93.3\% \\
    외국인 & 손실구간 $\times$ 환율   & $-0.0075$ & $0.0076$ & 86.7\% \\
    \addlinespace
    기관 & 모멘텀 $\times$ KOSPI200 & $-0.0001$ & $0.0005$ & 66.7\% \\
    기관 & 모멘텀 $\times$ 환율     & $+0.0003$ & $0.0006$ & 80.0\% \\
    기관 & 군집 $\times$ KOSPI200   & $-0.0004$ & $0.0006$ & 86.7\% \\
    기관 & 군집 $\times$ 환율       & $+0.0012$ & $0.0012$ & 93.3\% \\
    기관 & 손실구간 $\times$ KOSPI200 & $+0.0028$ & $0.0012$ & 97.8\% \\
    기관 & 손실구간 $\times$ 환율   & $+0.0007$ & $0.0008$ & 91.1\% \\
    \addlinespace
    개인 & 모멘텀 $\times$ KOSPI200 & $+0.0021$ & $0.0042$ & 80.0\% \\
    개인 & 모멘텀 $\times$ 환율     & $+0.0144$ & $0.0045$ & 100.0\% \\
    개인 & 군집 $\times$ KOSPI200   & $-0.0080$ & $0.0046$ & 100.0\% \\
    개인 & 군집 $\times$ 환율       & $+0.0092$ & $0.0081$ & 82.2\% \\
    개인 & 손실구간 $\times$ KOSPI200 & $-0.0022$ & $0.0046$ & 66.7\% \\
    개인 & 손실구간 $\times$ 환율   & $-0.0186$ & $0.0046$ & 100.0\% \\
    \bottomrule
  \end{tabular}
\end{table}

\begin{table}[htbp]
  \caption{도달한 해의 최적성 점검(45개 분할)}
  \label{tab:convergence}
  \small
  \centering
  \begin{tabular}{@{}lccc@{}}
    \toprule
    유형 & 훈련 MSE & 목적함수 격차(\%) & 코사인 \\
     & 감소율 & 평균 / 최대 & 평균 / 최소 \\
    \midrule
    외국인 & 14.1\% & $0.21/3.10$ & $0.994/0.948$ \\
    기관   & 1.5\%  & $0.17/0.48$ & $0.935/0.832$ \\
    개인   & 10.3\% & $0.92/1.40$ & $0.894/0.833$ \\
    \bottomrule
  \end{tabular}
\end{table}

%  필요 패키지: kotex, booktabs, graphicx, amsmath, array, url
%  그림 경로: ../figures_v2/  (scripts/make_paper_figures.py 산출물)
%  추가 참고문헌: paper-refs-additional.bib 병합 필요
%
%  [선택] 정리 환경 이름을 "Proposition"에서 "명제"로 바꾸려면 전처리부에
%  아래 세 줄을 추가한다. acmart가 이미 proposition을 정의하므로 먼저 해제한다.
%
%    \makeatletter
%    \let\proposition\relax \let\endproposition\relax
%    \makeatother
%    \theoremstyle{acmplain}
%    \newtheorem{proposition}{명제}[section]
%
%  추가하지 않으면 "Proposition 3.1"로 조판되며, 본문의 상호참조는 그대로
%  동작한다.
% =====================================================================

\section{방법}
\label{sec:method}

\subsection{설정: 대칭 추정}
\label{sec:setup}

투자자 유형 집합을 $\mathcal{I}=\{\text{외국인},\text{기관},\text{개인}\}$으로
둔다. 각 유형은 관측된 일별 수급이 해석 가능한 시장 상태 특징 위의 잠재 보상을
드러내는 하나의 에이전트다.

\paragraph{연속 행동}
유형 $i$의 $t$일 총매수금액과 총매도금액을 각각 $V^{buy}_{i,t}$,
$V^{sell}_{i,t}$라 하면 행동을
\begin{equation}
u_{i,t}=\frac{V^{buy}_{i,t}-V^{sell}_{i,t}}
             {V^{buy}_{i,t}+V^{sell}_{i,t}}\in[-1,1]
\label{eq:action}
\end{equation}
로 정의한다. 분모가 시장 전체 거래대금이 아니라 해당 유형의 총거래금액이므로
$u_{i,t}$는 그 유형 내부에서 매수와 매도 중 어느 쪽이 우세했는지를 나타내며,
유형별 거래 규모 차이에 영향받지 않는 공통 척도가 된다. 이진 매수ㆍ매도 라벨
대신 연속 강도를 사용하므로 이진화에서 버려지는 강도 정보가 보존된다.

\paragraph{대칭성}
세 유형은 동일한 상태 특징 $x_{i,t}\in\mathbb{R}^{3}$과 동일한 시장 컨텍스트
$C_t\in\mathbb{R}^{2}$, 그리고 동일한 모형 형태를 부여받는다. 유형별로 다른
특징을 미리 배정하지 않으며, 파라미터 $\theta_i=(\beta_i,B_i,\alpha_i)$만 각자의
데이터로 분리 학습한다. 따라서 결과에 나타나는 유형 간 차이는 설계가 심어둔 것이
아니라 데이터에서 학습된 것이며, 이 대칭성이 유형 간 계수의 직접 비교를
정당화한다.

\paragraph{시점 정렬}
학습 표본은 $(x_{i,t},C_t,a_{i,t+1})$이며 $a_{i,t+1}=u_{i,t+1}$이다. $t$일
종가와 수급으로 구성한 상태가 $t+1$일 행동을 설명하므로 미래 거래는 특징에
사용되지 않는다. 모멘텀과 손실구간 특징은 $t$일 정보를, 군집 특징은 $t-1$일
정보를 사용한다. 예측 대상이 $t+1$일 행동이므로 $t$일 정보도 이용 가능하며,
군집 특징의 추가 지연은 정보 제약이 아니라 동일 거래일 유형 간 수급의 회계적
상쇄를 예측 변수에서 배제하기 위한 설계 선택이다(그림~\ref{fig:crossinv}).

\subsection{보상 특징}
\label{sec:features}

세 특징은 모두 각 검증 분할의 학습 인덱스에서만 적합한 표준화 변환을 거친다.
따라서 이후의 모든 계수는 해당 특징이 학습 구간에서 1표준편차 변할 때의 효과로
읽는다.

\paragraph{(F1) 중기 모멘텀}
조정종가 또는 대체 종가 $P_t$로
\begin{equation}
x^{mom}_{i,t}=\log\!\left(\frac{P_t}{P_{t-20}}\right)
\label{eq:momentum}
\end{equation}
와 같이 계산한다. 20거래일은 약 1개월에 해당하는 사전 지정 창으로, 일별 가격
잡음을 완화하면서도 다음 날 행동과 연결되는 중기 추세를 표현한다. 자료에서
성능을 최적화해 고른 값이 아니라 해석 가능성을 위한 설계값이다. 양의 계수는
추세추종을, 음의 계수는 역추세를 뜻하며, 부호 기대는 유형별 모멘텀 거래
성향에 관한 선행 연구~\cite{grinblatt2000investment,choe1999foreign}에서
도출된다.

\paragraph{(F2) 군집행동}
유형 $i$를 제외한 두 유형 $j,k$의 전일 행동에 대한 횡단면 평균
\begin{equation}
x^{herd}_{i,t}=\frac{u_{j,t-1}+u_{k,t-1}}{2},\qquad j,k\neq i
\label{eq:herd}
\end{equation}
으로 정의한다. 시간평균이 아니라 같은 시점의 다른 두 유형에 대한 평균임에
유의한다. 20일 이동평균을 쓰면 즉각적 반응이 아니라 한 달 누적 수급 압력을
측정하게 되므로, 본 연구는 전일 반응에 초점을 맞춘다. 양의 계수는 다른 유형의
전일 방향을 따르는 행동을, 음의 계수는 반대 방향 행동을 뜻한다. 부호 기대는
기관 군집행동 문헌~\cite{lakonishok1992impact}에서 도출되나, 해당 문헌이 기관
\emph{내부}의 군집을 측정하는 반면 식~\eqref{eq:herd}은 \emph{유형 간} 관계를
측정한다는 점에서 대상이 다르다(\S\ref{sec:construct}).

\paragraph{(F3) 평균단가 대비 손실구간}
유형별 평균단가 대용치 $\bar P_{i,t}$와 당일 총거래 VWAP $P^{vwap}_{i,t}$로
\begin{equation}
x^{uw}_{i,t}=\max\!\left(
\frac{\bar P_{i,t}-P^{vwap}_{i,t}}{\bar P_{i,t}},\,0\right)
\label{eq:underwater}
\end{equation}
와 같이 정의한다. 0에서 절단하므로 이 변수는 거래가격이 평균단가 대용치보다
낮은 구간에서만 작동하며, 그 자체로 준거점 의존적 설계다. 평균단가는 총매수ㆍ
총매도 수량과 금액으로 유효 보유량을 재귀적으로 갱신해 계산하되, 전일 보유량에
망각계수 $\rho=0.98$을 곱해 오래된 거래의 영향을 줄이고 매도는 남은 물량의
평균단가를 바꾸지 않는다. $\rho=0.98$의 반감기는 약 34거래일이고 20거래일 후
약 66.8\%가 남으므로 최근 1--2개월의 거래기억을 강조하며, 이 값 역시 자료에서
최적화하지 않았다. 양의 계수는 평균단가 아래에서 다음 날 매수 강도가 커지는
방향이며, 부호 기대는 처분효과
문헌~\cite{shefrin1985disposition,odean1998investors}에서 도출된다.

이 대용치는 실제 계좌 보유량이 아니라 집계 거래로부터 구성한 근사치다. 같은
유형 안의 서로 다른 계좌가 수행한 매수와 매도가 넷팅되므로 ``대표 투자자의
평균단가''라는 개념 자체가 근사에 불과하며, 따라서 $x^{uw}$의 계수는
처분효과를 직접 검정한 결과로 해석할 수 없다.

\subsection{시장 컨텍스트}
\label{sec:context}

컨텍스트는 두 변수로 구성한다. (C1) KOSPI200 1일 수익률 $r_t$, (C2) 원/달러
환율 수준
\begin{equation}
z_t=\frac{\log e_t-\mu_{t-1,252}}{\sigma_{t-1,252}}
\label{eq:fxcontext}
\end{equation}
이다. 환율의 현재값은 사용하지만 평균과 표준편차는 $t-1$까지의 252거래일만
사용한다. 두 컨텍스트 역시 각 분할의 학습 인덱스에서만 다시 표준화하므로,
계수는 해당 컨텍스트가 학습 구간에서 1표준편차 변할 때의 효과로 읽는다.

컨텍스트와 보상 특징은 서로 독립이 아니다. 삼성전자가 KOSPI200에서 차지하는
비중 때문에 종목 모멘텀과 시장수익률이 겹치며, 특징--컨텍스트 상관은 $-0.29$에서
$0.22$ 범위에 있다(그림~\ref{fig:featcorr}). 따라서 개별 계수는 다른 항이
통제된 조건부 연관이며 단독으로 인과적 크기로 읽어서는 안 된다.

\subsection{근시안적 보상, 정책, 그리고 Lasso로의 환원}
\label{sec:model}

\paragraph{잠재 행동점수}
유형 $i$의 시점별 행동점수를
\begin{equation}
q_{i,t}=(\beta_i+B_iC_t)^{\top}x_{i,t}+\alpha_i^{\top}C_t
\label{eq:score}
\end{equation}
로 둔다. $\beta_i\in\mathbb{R}^{3}$은 컨텍스트가 평균일 때의 기본 보상 가중치,
$B_i\in\mathbb{R}^{3\times 2}$는 특징--컨텍스트 상호작용,
$\alpha_i\in\mathbb{R}^{2}$는 컨텍스트의 직접 주효과다. $\alpha$항은 종목
특징이 평균값일 때도 시장 국면이 행동 수준을 이동시키는 수준효과이고, $B$항은
국면에 따라 각 특징에 대한 민감도가 달라지는 기울기 효과다. 상호작용항을
포함하면서 하위차수 항 $C_t$를 누락하는 것은 표준적 사양 오류이므로,
$\alpha$의 포함은 자료로 정당화할 대상이 아니라 사양상 요구되는 것이다.

\paragraph{근시안적 오목 보상}
연속 행동 $a$에 대한 단기 보상을
\begin{equation}
R_i(x,C,a)=q_{i,t}\,a-\frac{\kappa}{2}a^{2}
\label{eq:reward}
\end{equation}
로 둔다. 이차항은 큰 순매수와 순매도에 대칭적 비용을 부과하고 보상을 행동에
대해 오목하게 만들어 내부해를 가능하게 한다. 이 항이 없으면 $[-1,1]$에서의
최대화 해가 $q$의 부호에 따라 거의 항상 $\pm1$이 되어 행동 \emph{강도}를
재현할 수 없다. 식~\eqref{eq:reward}을 $a$에 대해 최대화하면 닫힌형 정책
\begin{equation}
a^{*}_{i,t+1}=\mathrm{clip}\!\left(\frac{q_{i,t}}{\kappa},-1,1\right)
\label{eq:policy}
\end{equation}
을 얻는다. 곡률 $\kappa$는 식별을 위해 1로 고정한다. $\kappa=2$이면 같은 $q$에서
행동 크기가 $1/2$이 되지만 재학습 시 $q$의 계수를 2배로 바꾸면 같은 행동이
재현되므로, $\kappa=1$은 보상과 계수의 단위를 고정하는 정규화이지 곡률이
경험적으로 1이라는 주장이 아니다.

\paragraph{추정}
유형별 학습목적은 학습집합 $\mathcal{D}_i$에서
\begin{equation}
\min_{\theta_i}\ \frac{1}{|\mathcal{D}_i|}
\sum_{t\in\mathcal{D}_i}\bigl(a^{*}_{i,t+1}-a_{i,t+1}\bigr)^{2}
+\lambda_i\lVert\theta_i\rVert_1
\label{eq:loss}
\end{equation}
이며, Adam~\cite{kingma2015adam}으로 유형별 모형을 독립 학습한다.

\begin{proposition}[Lasso 환원]
\label{prop:lasso}
$z_{i,t}=[\,x_{i,t};\ \mathrm{vec}(x_{i,t}C_t^{\top});\ C_t\,]\in\mathbb{R}^{11}$,
$\theta_i=[\,\beta_i;\ \mathrm{vec}(B_i);\ \alpha_i\,]$로 두면
식~\eqref{eq:score}는 $q_{i,t}=\theta_i^{\top}z_{i,t}$이다. 학습 표본 전체에서
$|q_{i,t}|\le 1$이면 식~\eqref{eq:policy}의 clip이 작동하지 않아 $a^{*}=q$이고,
식~\eqref{eq:loss}은
\begin{equation}
\min_{\theta_i}\ \frac{1}{|\mathcal{D}_i|}
\sum_{t}\bigl(\theta_i^{\top}z_{i,t}-a_{i,t+1}\bigr)^{2}
+\lambda_i\lVert\theta_i\rVert_1
\label{eq:lasso}
\end{equation}
이 된다. 이는 다음 날 순매수 강도를 확장 상태특징 $z$에 회귀하는 $L_1$
회귀(Lasso)와 정확히 같다.
\end{proposition}

즉 $\kappa=1$의 근시안적 오목 보상 아래에서 ``선호를 복원한다''와 ``수급을
회귀한다''는 수치적으로 구별되지 않는다. 그림~\ref{fig:saturation}은
명제~\ref{prop:lasso}의 전제가 본 표본에서 실제로 성립함을 보인다. 세 유형 모두
잠재 행동점수의 분포가 행동 경계 $\pm1$로부터 크게 떨어져 있으며($|q|$의
최대값은 외국인 0.842, 개인 0.587, 기관 0.135), 45개 분할 전부에서 포화 비율이
정확히 0이다.

\paragraph{이 환원을 숨기지 않는 이유}
이 일치는 방법의 약점이 아니라 근시안 가정의 논리적 귀결이다. 에이전트가 한
걸음만 내다본다고 가정하면 오늘의 행동이 내일 이후에 영향을 주지 않으므로 내다볼
미래가 없고, 따라서 무엇을 원하는가(보상)와 무엇을 하는가(정책)가 하나로
겹친다. 최대엔트로피 IRL~\cite{ziebart2008maximum}의 확률적 정책
$\pi(a\mid x,C)\propto\exp(R/\tau)$를 식~\eqref{eq:reward}에 적용하면 $\pi$는
평균 $q$, 분산 $\tau$의 절단정규분포가 되어 로그가능도 최대화의 점추정치가
식~\eqref{eq:lasso}과 일치하므로, 이 환원은 특정 구현의 우연이 아니라
근시안ㆍ이차보상 설정 전반의 성질이다.

역강화학습 정식화의 구별되는 가치는 확장 경로에 있다. 상태 동역학과 다기간
가치를 도입하는 순간 보상은 정책과 갈라지고, 회귀로는 도달할 수 없는 구조적
복원으로 이어진다. 본 논문은 그 경로의 첫 단계이며, 성립 조건을 명시한 위에
해석과 검증을 더한다.

\subsection{검증 프로토콜}
\label{sec:protocol}

복원된 보상을 계수 부호와 크기의 보고로 끝내지 않고 다음 네 단계를 적용한다.

\begin{description}
\item[(V1) CPCV 표본외 안정성] 10개 시간 폴드 중 2개를 시험 폴드로 선택하는
조합형 정화 교차검증~\cite{lopez2018advances}을 사용한다. 조합 수는
$\binom{10}{2}=45$이며, 학습과 시험 사이에 1일 purge와 5일 embargo를 적용하고
모든 표준화 변환은 각 분할의 학습 인덱스에만 적합한다. 45개 분할에서 지배 부호가
나타난 분할의 비율(부호 일관성)을 보고하되, 분할들이 관측을 공유하므로 부호
일관성은 독립표본 $p$-값이 아니라 안정성 지표로만 해석한다.

\item[(V2) 달력월 블록 부트스트랩] 부호 일관성이 검정이 아니라는 한계를
보완하기 위해, 표본을 달력월 단위 블록으로 나누고 월 블록을 복원추출해 1{,}000회
재추정하여 계수의 95\% 신뢰구간을 구한다. 월 블록을 쓰는 이유는 일별 수급의
자기상관과 월중 패턴을 블록 내부에 보존하기 위해서다.

\item[(V3) Expanding walk-forward] CPCV는 시험 폴드가 학습 폴드보다 앞설 수
있으므로, 시간 방향을 엄격히 지키는 검증을 추가한다. 최초 400개 관측으로 학습해
5일 purge 후 다음 60거래일을 예측하고, 학습 구간을 60거래일씩 확장하며
반복한다(총 10개 창).

\item[(V4) 정규화 방식 교체] 결과가 $L_1$ 벌점의 산물이 아님을 확인하기 위해
동일 프로토콜을 $L_2$ 벌점과 무벌점으로 재실행하고 계수 부호를 비교한다.
\end{description}

\paragraph{평가지표}
행동 강도 오차는 MAE와 RMSE로, 방향 재현은
$\mathrm{sign}(a^{*})=\mathrm{sign}(a)$의 비율인 방향 정확도로 평가한다. 실제
행동과 예측 행동의 피어슨 상관, $|a^{*}|=1$인 포화 비율을 함께 보고한다. 유형 간
행동 표준편차가 다르므로 MAE와 RMSE의 절대값은 유형 간 비교에 쓰지 않고 각 유형
안에서 해석하며, 표준편차로 정규화한 오차를 병기한다. 참조 기준선으로 전일
자기행동을 그대로 예측하는 지속성 규칙을 함께 보고한다(\S\ref{sec:notclaim}).

\paragraph{구성 타당성 점검}
복원된 부호를 독립적으로 확립된 행동재무 규칙성에서 도출한 사전 기대와
대조한다(\S\ref{sec:construct}). 이 대조는 사후 해석이 아니라 사전에 정한 부호
기대와의 비교다.

% =====================================================================
\section{실험}
\label{sec:experiments}

\subsection{데이터와 설정}
\label{sec:data}

원자료는 삼성전자(005930) 일별 자료이며 가격, 거래대금, 유형별 총매수ㆍ총매도
금액과 수량, 원/달러 환율, KOSPI200 수익률을 포함한다. 252일 환율 이력, 20일
모멘텀, 다음 날 행동을 모두 계산할 수 있는 관측 중 최신 973행을 분석 표본으로
사용하였다. 상태일은 2022년 1월 6일--2025년 12월 29일이며, 마지막 상태는
2025년 12월 30일 행동을 목표로 한다.

단일 종목은 표본의 한계가 아니라 교차종목 교란 없이 유형 간 선호 이질성만
분리하기 위한 통제된 시연 설정이다. 종목 고정효과, 유동성 차이, 지수 편입 효과가
개입하지 않으므로 유형 간 계수 차이를 설계가 아닌 학습의 결과로 귀속할 수 있다.

\begin{table}[t]
  \caption{데이터 및 추정 설정}
  \label{tab:setup}
  \small
  \begin{tabular}{@{}>{\raggedright\arraybackslash}p{0.30\columnwidth}
                  >{\raggedright\arraybackslash}p{0.62\columnwidth}@{}}
    \toprule
    항목 & 설정 \\
    \midrule
    표본 & 973거래일, 2022-01-06--2025-12-29 \\
    유형 & 외국인, 기관, 개인 (대칭 추정) \\
    보상 특징 & 모멘텀(20일), 군집(전일), 손실구간($\rho=0.98$) \\
    컨텍스트 & KOSPI200 1일 수익률, USD/KRW 252일 수준 $z$ \\
    유형별 파라미터 & 11개 ($\beta$ 3, $B$ 6, $\alpha$ 2) \\
    공통 설정 & seed 42, CPU, $\kappa=1$, 학습구간 전용 표준화 \\
    외국인 & 75 epochs, batch 256, lr 0.001, $\lambda=0$ \\
    기관 & 10 epochs, batch 2048, lr 0.0005, $\lambda=0.01$ \\
    개인 & 20 epochs, batch 512, lr 0.001, $\lambda=0.0003$ \\
    검증 & CPCV 10 folds, test 2, purge 1, embargo 5, 45 splits \\
    안정성 & 월 블록 부트스트랩 1{,}000회, walk-forward 10창, 정규화 교체 \\
    \bottomrule
  \end{tabular}
\end{table}

행동 평균은 외국인 $-0.0063$, 기관 $-0.0100$, 개인 $0.0035$로 세 유형 모두 0에
가까우나 표준편차는 개인(0.334)이 가장 크고 기관(0.125)이 가장 작다. 1차
자기상관은 외국인 0.401, 기관 0.149, 개인 0.345다.

동일 거래일 유형 간 행동은 강한 음의 관계를 갖는다(그림~\ref{fig:crossinv}).
외국인--개인 상관이 $-0.772$, 기관--개인이 $-0.499$이며, 이는 한 유형의 매수가
다른 유형의 매도로 나타날 수밖에 없는 시장 청산 제약을 반영한다. 이 사실은
\S\ref{sec:beta}과 \S\ref{sec:alpha}의 해석 범위를 제한하는 근거로 다시
사용된다.

\subsection{표본외 행동 재현}
\label{sec:oos}

\begin{figure}[t]
  \centering
  \includegraphics[width=\columnwidth]{../figures_v2/fig2_actual_vs_prediction.pdf}
  \caption{투자자별 실제 행동과 CPCV 표본외 예측 행동. 각 날짜의 9개 CPCV
  시험예측을 먼저 평균한 뒤 실제값과 예측값에 20거래일 이동평균을 적용했다.
  괄호 안 상관은 이동평균 계열 간 값이며 표~\ref{tab:performance}의 일별
  표본외 상관과 직접 비교하지 않는다.}
  \Description{외국인, 기관, 개인의 실제 순매수 행동과 모형 예측을 시간축 위에
  겹쳐 그린 세 개의 선그래프. 기관 패널에서 예측선이 0 근방에 평탄하게 놓여 있다.}
  \label{fig:prediction}
\end{figure}

그림~\ref{fig:prediction}은 세 유형의 실제 순매수 행동과 표본외 예측 행동을
나란히 보인다. 세 가지가 한눈에 읽힌다. 첫째, 외국인과 개인에서는 예측 계열이
실제 계열의 국면 전환을 따라간다(이동평균 계열 간 상관 0.69, 0.76). 둘째, 예측
진폭이 실제보다 뚜렷하게 작으며, 이는 $L_1$ 정규화에 의한 계수 축소의 결과다.
셋째, 기관의 예측 계열은 사실상 0선에 붙어 있다(0.02). \S\ref{sec:null}에서
보고하는 null이 계수 표가 아니라 예측 시계열에서 먼저 드러난다.

\begin{table}[t]
  \caption{CPCV 표본외 성능(45개 분할 평균 $\pm$ 표준편차)과 지속성 기준선}
  \label{tab:performance}
  \small
  \centering
  \begin{tabular}{@{}lccc@{}}
    \toprule
    지표 & 외국인 & 기관 & 개인 \\
    \midrule
    \multicolumn{4}{@{}l}{\emph{주모형}}\\
    방향 정확도   & $0.636\pm0.039$ & $0.542\pm0.038$ & $0.608\pm0.044$ \\
    상관계수      & $0.278\pm0.110$ & $0.068\pm0.064$ & $0.245\pm0.118$ \\
    MAE           & $0.195\pm0.019$ & $0.094\pm0.008$ & $0.264\pm0.029$ \\
    RMSE          & $0.246\pm0.021$ & $0.125\pm0.012$ & $0.318\pm0.030$ \\
    MAE / 행동 SD & $0.759$ & $0.756$ & $0.791$ \\
    포화 비율     & $0.000$ & $0.000$ & $0.000$ \\
    \addlinespace
    \multicolumn{4}{@{}l}{\emph{지속성 기준선} ($a_{t+1}\leftarrow a_t$)}\\
    방향 정확도   & $\mathbf{0.642}$ & $\mathbf{0.547}$ & $\mathbf{0.616}$ \\
    상관계수      & $\mathbf{0.359}$ & $\mathbf{0.136}$ & $\mathbf{0.311}$ \\
    MAE           & $0.223$ & $0.123$ & $0.299$ \\
    \addlinespace
    행동의 1차 자기상관 & $0.401$ & $0.149$ & $0.345$ \\
    \bottomrule
  \end{tabular}
\end{table}

\begin{figure}[t]
  \centering
  \includegraphics[width=\columnwidth]{../figures_v2/fig5_oos_performance.pdf}
  \caption{표본외 성능의 분할 간 분포와 자명한 기준선. 실선은 주모형 평균,
  파선은 지속성 기준선이다. 세 유형 모두에서 파선이 실선보다 위에 있다.}
  \Description{방향 정확도와 상관계수에 대해 세 투자자 유형의 45개 분할 값을
  바이올린 및 점으로 그린 두 개의 패널.}
  \label{fig:oos}
\end{figure}

표~\ref{tab:performance}과 그림~\ref{fig:oos}에서 평균만 보면 가려지는 두 가지가
드러난다. 첫째, 분할 간 산포가 크다. 특히 상관계수는 외국인이 $-0.05$에서
$0.45$까지 걸쳐 있어 단일 분할의 수치로 유형을 비교하는 것은 의미가 없다. 둘째,
자명한 지속성 기준선이 방향 정확도와 상관계수에서 주모형을 앞선다. $a_t$는 $t$일
종가 이후 공개되므로 이는 정보상 이용 가능한 규칙이다. 주모형이 앞서는 것은
MAE뿐인데, 이는 정규화로 예측 진폭이 축소되어 크기 오차가 줄어든 결과이므로
정보량 우위로 읽을 수 없다.

이 격차의 원인은 명확하다. 행동의 1차 자기상관이 $0.401/0.149/0.345$인데 본
연구의 보상 특징 집합은 자기 지연행동을 포함하지 않는다. 특징을 (F1)--(F3)로
제한한 것은 세 유형이 독립적으로 확립된 행동재무 규칙성에 대응하는 특징만
공유하도록 설계했기 때문이며, 예측 성능을 극대화하려는 설계가 아니다. 본 연구는
이 결과를 은폐하지 않고 표~\ref{tab:performance}에 함께 싣는다.
표~\ref{tab:performance}과 그림~\ref{fig:oos}는 모형이 잡음만 학습한 것이 아님을
보이는 용도로 사용되며, 예측 성능의 우위를 주장하는 근거로 사용되지 않는다.

모든 유형ㆍ모든 분할에서 포화 비율이 정확히 0이므로 명제~\ref{prop:lasso}의 전제
$|q|\le1$이 표본 전체에서 성립한다. 즉 본 추정은 실제로 식~\eqref{eq:lasso}의
Lasso와 동일하며, 성능이 $\pm1$ 극단 예측에 의존하지도 않는다.

\subsection{복원된 보상: 모멘텀 축의 거울상}
\label{sec:beta}

\begin{figure*}[t]
  \centering
  \includegraphics[width=\textwidth]{../figures_v2/fig1_reward_weights_bar.pdf}
  \caption{투자자별로 복원된 기본 보상가중치 $\beta$. 파란색은 양, 붉은색은 음의
  평균 계수이며, 오차막대는 45개 분할의 1 표준편차, 괄호 안은 부호 일관성이다.}
  \Description{외국인, 기관, 개인 각각에 대해 모멘텀, 군집, 손실구간 계수를 그린
  세 개의 막대그래프. 외국인 모멘텀은 양, 개인 모멘텀은 음이며 기관의 세 막대는
  거의 보이지 않는다.}
  \label{fig:betabar}
\end{figure*}

\begin{figure*}[t]
  \centering
  \includegraphics[width=\textwidth]{../figures_v2/fig3_beta_alpha_forest.pdf}
  \caption{보상가중치와 컨텍스트 주효과의 분할 간 안정성. 회색 점은 개별 CPCV
  분할, 굵은 오차선은 45개 분할의 표준편차, 옅은 가로선은 달력월 블록 부트스트랩
  95\% 신뢰구간이다. 속이 빈 표식은 신뢰구간이 0을 포함하는 항이다.}
  \Description{패널 A는 기본 보상가중치, 패널 B는 컨텍스트 주효과에 대한 forest
  plot으로, 각 항의 분할별 추정치와 신뢰구간을 보인다.}
  \label{fig:forest}
\end{figure*}

그림~\ref{fig:betabar}는 세 유형에 대해 복원된 기본 보상가중치를 나란히 보인다.
동일한 특징 집합과 동일한 모형 형태를 부여하고 파라미터만 분리 학습했으므로, 세
패널의 높이 차이는 전적으로 데이터에서 학습된 것이다.
그림~\ref{fig:forest}는 같은 $\beta$와 컨텍스트 주효과 $\alpha$를 분할 수준까지
펼쳐 안정성을 보인다.

\paragraph{모멘텀 축의 거울상}
외국인의 모멘텀 가중치는 $+0.0496$으로 45개 분할 전부에서 양수였고, 개인은
$-0.0304$로 45개 분할 전부에서 음수였다. 그림~\ref{fig:betabar}에서 두 막대가
0선을 사이에 두고 반대 방향을 향하며, 그림~\ref{fig:forest}에서 두 부트스트랩
신뢰구간은 서로 겹치지 않고 각각 0을 제외한다. 동일한 특징 집합을 유형별로 독립
추정했는데도 두 유형이 모멘텀에서 정반대로 갈린다. 외국인은 추세추종, 개인은
역추세로 나타나며, 이는 선행 연구~\cite{grinblatt2000investment,choe1999foreign}의
부호 기대와 일치한다.

\paragraph{군집}
군집 계수는 세 유형 모두 45개 분할 전부에서 음수였고, 외국인($-0.0402$)과
개인($-0.0322$)은 부트스트랩 신뢰구간도 0을 제외한다. 그러나 이 결과는 행동적
해석과 시장 청산 제약을 분리하지 못한다. 그림~\ref{fig:crossinv}가 보이듯 동일
거래일 유형 간 행동 상관이 $-0.77$에 이르므로, 음의 군집 계수 중 상당 부분이
회계적 제약에서 발생할 수 있다. 식~\eqref{eq:herd}가 1거래일 지연값을 쓰므로
동시점 제약이 직접 적용되지는 않으나, 행동의 자기상관을 통해 제약이 전이될 수
있다. 따라서 본 연구는 이를 ``관측된 집계 수급이 유형 간 음의 교차 연관을
보인다''는 서술적 결과로 제한하고, 각 유형이 독립적으로 반대매매 성향을 갖는다는
행동적 주장으로 확장하지 않는다.

\paragraph{손실구간}
개인의 손실구간 가중치는 $+0.0293$으로 45개 분할 전부에서 양수였다. 평균단가
대용치 아래에서 다음 날 매수 강도가 커지는 방향이며, 처분효과
문헌~\cite{shefrin1985disposition,odean1998investors}의 기준가격 의존성과
정합한다. 외국인도 $+0.0094$(95.6\%)로 같은 부호이나 크기가 개인의 $1/3$
수준이다.

다만 그림~\ref{fig:forest}에서 보듯 손실구간은 세 유형 모두 부트스트랩 신뢰구간이
0을 포함한다. 개인의 구간은 $[-0.0122,+0.0665]$로, 45개 분할에서 부호가 100\%
일관됨에도 불구하고 월 블록 재표집에서는 부호가 뒤집히는 경우가 5.5\% 발생한다.
이는 CPCV 분할들이 관측을 공유하기 때문에 부호 일관성이 불확실성을 과소평가한다는
\S\ref{sec:protocol}의 경고가 실제로 작동한 사례다. 따라서 본 연구는 처분효과와
정합하는 부호가 관측되었다고 보고하되, 이를 통계적으로 확립된 결과로 주장하지
않는다.

\subsection{컨텍스트 반응: 거울상의 재현}
\label{sec:alpha}

모멘텀 축에서 관찰된 거울상이 컨텍스트 주효과에서도 그대로
재현된다(그림~\ref{fig:forest}의 패널 B). 외국인은 시장 상승일에 다음 날 순매수
방향이 강해지고($+0.0361$) 원화 약세 국면에서는 순매도 방향으로
이동하는($-0.0261$) 반면, 개인은 두 컨텍스트 모두에서 정확히
반대였다($-0.0246$, $+0.0298$). 네 계수 모두 45개 분할 전부에서 부호가
유지되었고 부트스트랩 신뢰구간도 모두 0을 제외한다. 기관의 두 계수는 부호
일관성이 95\% 이상이나 절대 크기가 외국인ㆍ개인의 $1/8$ 수준이고 신뢰구간이 0을
포함한다.

이 거울상 역시 군집 계수와 같은 해석상 유보가 필요하다.
그림~\ref{fig:crossinv}의 $-0.772$라는 상관 하에서 두 유형의 컨텍스트 반응이
반대 부호로 나타나는 것은 부분적으로 시장 청산 제약을 반영할 수 있다. 본 논문은
이를 두 유형이 시장ㆍ환율 국면에서 반대편 수급을 형성한다는 서술적 사실로
제시한다.

\paragraph{상호작용}
12개 상호작용 계수 중 45개 분할 전부에서 부호가 유지된 것은 외국인의
모멘텀 $\times$ KOSPI200($-0.0139$), 개인의 모멘텀 $\times$ 환율($+0.0144$),
개인의 군집 $\times$ KOSPI200($-0.0080$), 개인의 손실구간 $\times$
환율($-0.0186$) 네 개이며, 이 중 부트스트랩을 통과하는 것은 앞의 하나와 마지막 하나다. 12개 계수를
동시에 검토했으므로 본 연구는 상호작용을 탐색적 결과로만 제시한다.

\subsection{안정성 프로토콜의 결과}
\label{sec:stability}

\paragraph{(V2) 달력월 블록 부트스트랩}
그림~\ref{fig:forest}의 옅은 가로선이 그 결과다. 부트스트랩은 CPCV 부호
일관성보다 엄격한 기준이며, 결과는 셋으로 갈린다. 첫째, 모멘텀 거울상과 컨텍스트
거울상은 부트스트랩에서도 유지된다. 외국인ㆍ개인의 모멘텀, 군집, 두 컨텍스트
주효과 등 8개 항이 모두 0을 제외한다. 둘째, 손실구간 가중치는 통과하지 못한다.
셋째, 기관은 어떤 항도 0을 제외하지 못하며, 그림~\ref{fig:forest}에서 기관의
여섯 표식이 모두 속이 비어 있다.

\begin{figure*}[t]
  \centering
  \includegraphics[width=\textwidth]{../figures_v2/fig6_walk_forward.pdf}
  \caption{Expanding walk-forward 10개 창의 계수 경로와 방향 정확도. 성능은
  창마다 흔들리지만 계수 부호는 유지된다.}
  \Description{확장 창 번호를 가로축으로 하여 세 유형의 모멘텀 계수, 두 컨텍스트
  주효과, 방향 정확도를 그린 네 개의 선그래프.}
  \label{fig:walkforward}
\end{figure*}

\paragraph{(V3) Expanding walk-forward}
그림~\ref{fig:walkforward}가 그 결과다. 두 가지가 동시에 읽힌다. 성능은 창마다
크게 흔들린다. 외국인 방향 정확도는 0.43에서 0.85 사이를 오가며 평균
$0.610\pm0.116$, 상관계수는 $0.131\pm0.162$로 CPCV의 0.278에서 절반 이하로
떨어진다. 개인도 $0.245\rightarrow0.129$다. CPCV는 시험 폴드를 학습 폴드보다
앞세울 수 있는 반면 walk-forward는 그렇지 않으므로, 본 연구는 이 저하를 시간
구조에서 오는 일반화 한계로 그대로 보고한다.

반면 계수 경로는 흔들리지 않는다. 외국인의 모멘텀 계수는 10개 창 전부에서 양수,
개인은 전부에서 음수이며 두 경로가 한 번도 교차하지 않는다. 두 컨텍스트
주효과에서도 같은 거울상이 창 전체에 걸쳐 유지된다. 기관의 세 경로는 모두 0선에
붙어 있다. 즉 성능은 시간에 따라 흔들리지만 복원된 선호의 방향은 흔들리지
않는다.

\begin{figure*}[t]
  \centering
  \includegraphics[width=\textwidth]{../figures_v2/fig7_institution_null.pdf}
  \caption{기관의 수렴된 null. 좌: 복원된 선호의 크기(로그 축). 우: 정규화 방식
  교체에 대한 민감도이며 세 패널은 각각 독립 $y$축을 쓴다.}
  \Description{왼쪽은 유형별 계수 절대값을 로그 축 막대로, 오른쪽은 L1, L2,
  무벌점 하의 계수를 유형별로 그린 세 개의 선그래프.}
  \label{fig:null}
\end{figure*}

\paragraph{(V4) 정규화 방식 교체}
그림~\ref{fig:null}이 그 결과다. 외국인과 개인의 계수는 벌점 방식에 사실상
불변이다(소수점 넷째 자리까지 일치). 따라서 모멘텀 거울상은 $L_1$ 벌점의 산물이
아니다. 반면 기관의 계수만 벌점에 따라 이동하며(군집
$-0.0076\rightarrow-0.0124$, 손실구간 $-0.0000\rightarrow-0.0018$), 그 $y$축
범위는 다른 두 유형의 $1/20$ 수준이다. 계수가 벌점 선택에 좌우된다는 것은
데이터가 그 값을 결정하지 못한다는 뜻이다.

\subsection{기관: 수렴된 null}
\label{sec:null}

기관 수급에 대해 본 연구는 일별 빈도에서 안정적으로 식별되는 선호가 없다는
결과를 얻었고, 이를 은폐하지 않고 발견으로 보고한다. 근거는 다섯이며 모두 그림에
나타난다.

\begin{enumerate}
\item \textbf{예측의 퇴화.} 그림~\ref{fig:prediction} 중간 패널에서 기관의 예측
계열은 실제 계열이 $\pm0.1$ 사이를 오가는 동안 0선에 붙어 있다.
\item \textbf{크기.} 그림~\ref{fig:betabar}에서 기관 패널의 세 막대는 다른 두
유형과 같은 $y$축 위에서 거의 보이지 않는다. 군집을 제외하면 모멘텀 $-0.0006$,
손실구간 $-0.0008$로 $1/40$ 수준이다.
\item \textbf{불확실성.} 그림~\ref{fig:forest}에서 기관의 여섯 표식은 모두 속이
비어 있다.
\item \textbf{벌점 의존성.} 그림~\ref{fig:null} 우측 중앙 패널에서 기관만 벌점에
따라 계수가 이동한다.
\item \textbf{시간 불안정.} 그림~\ref{fig:walkforward}에서 기관의 경로는 0선
근방에 붙어 있으며, 모멘텀 부호 일관성은 10개 창 중 40\%, 손실구간은 0\%다.
\end{enumerate}

\begin{figure*}[t]
  \centering
  \includegraphics[width=\textwidth]{../figures_v2/fig4_training_convergence.pdf}
  \caption{45개 CPCV 분할에서의 학습 수렴. 옅은 회색선은 개별 분할, 굵은 선과
  띠는 평균 $\pm$ 1 표준편차다. 주석은 도달한 해가 동일 목적함수의 폐쇄형
  최적해 대비 남긴 상대 격차다.}
  \Description{에포크를 가로축으로 하여 세 유형의 훈련 MSE 곡선을 그린 세 개의
  패널. 기관 패널의 감소폭이 가장 작다.}
  \label{fig:convergence}
\end{figure*}

이 null은 최적화 실패가 아니다. 그림~\ref{fig:convergence}는 45개 분할의 학습
곡선을 유형별로 보인다. 기관의 훈련 MSE는 학습 구간 전체에서 1.5\%만 감소하며
곡선이 시각적으로 평탄해 보이지 않는다. 그러나 이는 학습 부족이 아니라 애초에
줄일 수 있는 오차가 거의 없기 때문이다. 명제~\ref{prop:lasso}에 따라
식~\eqref{eq:loss}은 볼록 문제이므로 전역 최적해가 유일하며, 도달한 해와 동일
목적함수의 폐쇄형 최적해를 비교하면 기관의 목적함수 상대 격차는 평균 0.17\%,
최대 0.48\%에 불과하다(외국인 0.21\%, 개인 0.92\%). 즉 세 유형 모두 도달 가능한
최적해에 근접했으며, 기관에서만 그 최적해의 계수가 0 근방이다. 데이터가 기관에
대해 결정하는 선호가 없는 것이지, 최적화가 그것을 찾지 못한 것이 아니다.

경제적으로 이 null은 자연스럽게 읽힌다. ``기관''은 연기금, 투신, 보험, 사모,
금융투자를 합산한 범주이므로 서로 다른 목적함수를 가진 주체들의 수급이 상쇄된다.
유일하게 일관된 음의 군집 계수는 이들이 공통적으로 다른 유형의 반대편에 서는
유동성 공급자 역할을 한다는 해석과 정합하나, 부트스트랩 신뢰구간이 경계에
걸치므로 확정하지 않는다.

\subsection{구성 타당성 점검}
\label{sec:construct}

\begin{table}[t]
  \caption{사전 부호 기대와 복원 결과}
  \label{tab:construct}
  \small
  \begin{tabular}{@{}>{\raggedright\arraybackslash}p{0.26\columnwidth}
                  >{\raggedright\arraybackslash}p{0.15\columnwidth}
                  >{\raggedright\arraybackslash}p{0.22\columnwidth}
                  >{\raggedright\arraybackslash}p{0.24\columnwidth}@{}}
    \toprule
    규칙성 & 사전 기대 & 복원 결과 & 판정 \\
    \midrule
    외국인의 추세거래~\cite{grinblatt2000investment,choe1999foreign}
      & $\beta^{mom}>0$ & $+0.0496$ (100\%, CI 0 제외) & 부합 \\
    개인의 역추세 거래~\cite{grinblatt2000investment,kaniel2008individual}
      & $\beta^{mom}<0$ & $-0.0304$ (100\%, CI 0 제외) & 부합 \\
    처분효과~\cite{shefrin1985disposition,odean1998investors}
      & $\beta^{uw}>0$ & $+0.0293$ (100\%, CI 0 포함) & 부분 부합 \\
    기관의 군집행동~\cite{lakonishok1992impact}
      & $\beta^{herd}>0$ & $-0.0044$ (100\%, CI 0 포함) & 불일치 \\
    \bottomrule
  \end{tabular}
\end{table}

표~\ref{tab:construct}이 대조 결과다. 세 유형에 동일한 특징과 동일한 모형 형태를
부여하고 파라미터만 분리 학습했음에도, 복원된 부호가 각각 독립적으로 확립된
규칙성과 부합하고 두 유형이 모멘텀에서 정반대로 갈린다. 이는 복원된 보상이 임의의
적합 결과가 아니라는 구성 타당성의 증거다.

불일치 항목은 측정 대상이 다르기 때문이다. Lakonishok
외~\cite{lakonishok1992impact}는 기관 \emph{내부} 주체들이 같은 방향으로 몰리는
정도를 측정하는 반면, 식~\eqref{eq:herd}은 기관이 \emph{다른 두 유형}의 전일
수급과 같은 방향으로 움직이는지를 측정한다. 따라서 본 결과는 해당 문헌의 반증이
아니며, 유형 내부 군집을 검정하려면 기관 세부 주체별 자료가 필요하다.

\subsection{본 연구가 주장하지 않는 것}
\label{sec:notclaim}

프로토콜의 마지막 단계로, 본 결과가 지지하지 않는 주장을 명시한다.

\begin{enumerate}
\item \textbf{수익률 예측가능성.} 본 연구의 지표는 수급의 방향과 강도에 대한
재현 오차이며 수익률, 샤프지수, 거래비용을 포함하지 않는다.
표~\ref{tab:performance}의 수치로부터 어떤 매매 규칙도 도출되지 않는다.

\item \textbf{수급 예측 성능의 우위.} 표~\ref{tab:performance}과
그림~\ref{fig:oos}에서 보고했듯 지속성 기준선이 방향 정확도와 상관계수에서
주모형을 앞선다. 본 연구의 목적은 예측 성능이 아니라 해석 가능한 공통 척도에서의
선호 복원이며, 두 목적은 다른 특징 집합을 요구한다. 예비 점검에서 지연
자기행동을 네 번째 특징으로 추가하면 모형이 기준선을 넘어서는 반면(방향 정확도
$0.646/0.569/0.629$), 개인의 $\alpha^{KOSPI}$ 부호가 반전되고 손실구간 계수가
소멸한다. 즉 본 논문이 보고한 거울상 중 일부는 수급 자기상관과 분리되지 않는다.

\item \textbf{인과효과.} 복원된 가중치는 표준화된 관측 특징 사이의 조건부
연관이다. 그림~\ref{fig:featcorr}이 보이듯 컨텍스트와 보상 특징 사이에 $-0.29$에서
$0.22$ 범위의 상관이 존재하므로 개별 계수의 크기는 다른 항의 통제 조건에
의존한다.

\item \textbf{유형 간 음의 연관의 행동적 해석.} 군집 계수와 거울상 컨텍스트
반응은 시장 청산 제약과 분리되지 않는다(그림~\ref{fig:crossinv}).

\item \textbf{처분효과의 검정.} 개인의 손실구간 부호는 문헌과 정합하나 월 블록
부트스트랩을 통과하지 못했고(그림~\ref{fig:forest}), 평균단가 대용치는 계좌별
실현손익과 보유기간이 아니라 집계 거래에서 구성한 근사치다.

\item \textbf{종목 간 일반화.} 단일 종목 설정은 유형 간 이질성 분리를 위한
통제이며, 다른 종목과 시장으로의 확장은 검증되지 않았다.
\end{enumerate}

% =====================================================================
\appendix

\section{보충 자료}
\label{sec:appendix}

본문에서 참조한 보조 그림과 표를 모은다. 그림~\ref{fig:saturation}은
명제~\ref{prop:lasso}의 전제를, 그림~\ref{fig:featcorr}은 특징과 컨텍스트의
상관을, 그림~\ref{fig:crossinv}은 유형 간 동시점 행동을 보인다.
표~\ref{tab:interactions}은 본문에서 탐색적 결과로만 제시한 상호작용 계수
전체이고, 표~\ref{tab:convergence}은 도달한 해의 최적성 점검이다.

\begin{figure}[htbp]
  \centering
  \includegraphics[width=\columnwidth]{../figures_v2/figA1_no_saturation.pdf}
  \caption{잠재 행동점수의 분포와 행동 경계. $|q|\le1$이므로
  명제~\ref{prop:lasso}의 전제가 표본 전체에서 성립하며, 45개 분할의 포화 비율은
  정확히 0이다.}
  \Description{세 유형의 잠재 행동점수 히스토그램과 $\pm1$ 경계선.}
  \label{fig:saturation}
\end{figure}

\begin{figure}[htbp]
  \centering
  \includegraphics[width=\columnwidth]{../figures_v2/figA2_feature_correlation.pdf}
  \caption{보상 특징과 컨텍스트의 상관행렬(외국인 특징 기준, 973개 관측).}
  \Description{다섯 변수 사이의 상관계수를 색과 숫자로 표시한 히트맵.}
  \label{fig:featcorr}
\end{figure}

\begin{figure}[htbp]
  \centering
  \includegraphics[width=\columnwidth]{../figures_v2/figA3_cross_investor_actions.pdf}
  \caption{동일 거래일 유형 간 행동. 점선은 완전 상쇄선이며, 외국인과 개인의
  행동이 그 근방에 밀집한다.}
  \Description{유형 쌍별 행동 산점도 세 개와 각 쌍의 상관계수.}
  \label{fig:crossinv}
\end{figure}

\begin{table}[htbp]
  \caption{상호작용 계수 $B$(45개 분할)}
  \label{tab:interactions}
  \small
  \centering
  \begin{tabular}{@{}llrrr@{}}
    \toprule
    유형 & 특징 $\times$ 컨텍스트 & 평균 & 표준편차 & 부호 일관성 \\
    \midrule
    외국인 & 모멘텀 $\times$ KOSPI200 & $-0.0139$ & $0.0052$ & 100.0\% \\
    외국인 & 모멘텀 $\times$ 환율     & $-0.0030$ & $0.0057$ & 71.1\% \\
    외국인 & 군집 $\times$ KOSPI200   & $+0.0016$ & $0.0074$ & 66.7\% \\
    외국인 & 군집 $\times$ 환율       & $+0.0037$ & $0.0073$ & 68.9\% \\
    외국인 & 손실구간 $\times$ KOSPI200 & $-0.0062$ & $0.0040$ & 93.3\% \\
    외국인 & 손실구간 $\times$ 환율   & $-0.0075$ & $0.0076$ & 86.7\% \\
    \addlinespace
    기관 & 모멘텀 $\times$ KOSPI200 & $-0.0001$ & $0.0005$ & 66.7\% \\
    기관 & 모멘텀 $\times$ 환율     & $+0.0003$ & $0.0006$ & 80.0\% \\
    기관 & 군집 $\times$ KOSPI200   & $-0.0004$ & $0.0006$ & 86.7\% \\
    기관 & 군집 $\times$ 환율       & $+0.0012$ & $0.0012$ & 93.3\% \\
    기관 & 손실구간 $\times$ KOSPI200 & $+0.0028$ & $0.0012$ & 97.8\% \\
    기관 & 손실구간 $\times$ 환율   & $+0.0007$ & $0.0008$ & 91.1\% \\
    \addlinespace
    개인 & 모멘텀 $\times$ KOSPI200 & $+0.0021$ & $0.0042$ & 80.0\% \\
    개인 & 모멘텀 $\times$ 환율     & $+0.0144$ & $0.0045$ & 100.0\% \\
    개인 & 군집 $\times$ KOSPI200   & $-0.0080$ & $0.0046$ & 100.0\% \\
    개인 & 군집 $\times$ 환율       & $+0.0092$ & $0.0081$ & 82.2\% \\
    개인 & 손실구간 $\times$ KOSPI200 & $-0.0022$ & $0.0046$ & 66.7\% \\
    개인 & 손실구간 $\times$ 환율   & $-0.0186$ & $0.0046$ & 100.0\% \\
    \bottomrule
  \end{tabular}
\end{table}

\begin{table}[htbp]
  \caption{도달한 해의 최적성 점검(45개 분할)}
  \label{tab:convergence}
  \small
  \centering
  \begin{tabular}{@{}lccc@{}}
    \toprule
    유형 & 훈련 MSE & 목적함수 격차(\%) & 코사인 \\
     & 감소율 & 평균 / 최대 & 평균 / 최소 \\
    \midrule
    외국인 & 14.1\% & $0.21/3.10$ & $0.994/0.948$ \\
    기관   & 1.5\%  & $0.17/0.48$ & $0.935/0.832$ \\
    개인   & 10.3\% & $0.92/1.40$ & $0.894/0.833$ \\
    \bottomrule
  \end{tabular}
\end{table}

%  필요 패키지: kotex, booktabs, graphicx, amsmath, array, url
%  그림 경로: ../figures_v2/  (scripts/make_paper_figures.py 산출물)
%  추가 참고문헌: paper-refs-additional.bib 병합 필요
%
%  [선택] 정리 환경 이름을 "Proposition"에서 "명제"로 바꾸려면 전처리부에
%  아래 세 줄을 추가한다. acmart가 이미 proposition을 정의하므로 먼저 해제한다.
%
%    \makeatletter
%    \let\proposition\relax \let\endproposition\relax
%    \makeatother
%    \theoremstyle{acmplain}
%    \newtheorem{proposition}{명제}[section]
%
%  추가하지 않으면 "Proposition 3.1"로 조판되며, 본문의 상호참조는 그대로
%  동작한다.
% =====================================================================

\section{방법}
\label{sec:method}

\subsection{설정: 대칭 추정}
\label{sec:setup}

투자자 유형 집합을 $\mathcal{I}=\{\text{외국인},\text{기관},\text{개인}\}$으로
둔다. 각 유형은 관측된 일별 수급이 해석 가능한 시장 상태 특징 위의 잠재 보상을
드러내는 하나의 에이전트다.

\paragraph{연속 행동}
유형 $i$의 $t$일 총매수금액과 총매도금액을 각각 $V^{buy}_{i,t}$,
$V^{sell}_{i,t}$라 하면 행동을
\begin{equation}
u_{i,t}=\frac{V^{buy}_{i,t}-V^{sell}_{i,t}}
             {V^{buy}_{i,t}+V^{sell}_{i,t}}\in[-1,1]
\label{eq:action}
\end{equation}
로 정의한다. 분모가 시장 전체 거래대금이 아니라 해당 유형의 총거래금액이므로
$u_{i,t}$는 그 유형 내부에서 매수와 매도 중 어느 쪽이 우세했는지를 나타내며,
유형별 거래 규모 차이에 영향받지 않는 공통 척도가 된다. 이진 매수ㆍ매도 라벨
대신 연속 강도를 사용하므로 이진화에서 버려지는 강도 정보가 보존된다.

\paragraph{대칭성}
세 유형은 동일한 상태 특징 $x_{i,t}\in\mathbb{R}^{3}$과 동일한 시장 컨텍스트
$C_t\in\mathbb{R}^{2}$, 그리고 동일한 모형 형태를 부여받는다. 유형별로 다른
특징을 미리 배정하지 않으며, 파라미터 $\theta_i=(\beta_i,B_i,\alpha_i)$만 각자의
데이터로 분리 학습한다. 따라서 결과에 나타나는 유형 간 차이는 설계가 심어둔 것이
아니라 데이터에서 학습된 것이며, 이 대칭성이 유형 간 계수의 직접 비교를
정당화한다.

\paragraph{시점 정렬}
학습 표본은 $(x_{i,t},C_t,a_{i,t+1})$이며 $a_{i,t+1}=u_{i,t+1}$이다. $t$일
종가와 수급으로 구성한 상태가 $t+1$일 행동을 설명하므로 미래 거래는 특징에
사용되지 않는다. 모멘텀과 손실구간 특징은 $t$일 정보를, 군집 특징은 $t-1$일
정보를 사용한다. 예측 대상이 $t+1$일 행동이므로 $t$일 정보도 이용 가능하며,
군집 특징의 추가 지연은 정보 제약이 아니라 동일 거래일 유형 간 수급의 회계적
상쇄를 예측 변수에서 배제하기 위한 설계 선택이다(그림~\ref{fig:crossinv}).

\subsection{보상 특징}
\label{sec:features}

세 특징은 모두 각 검증 분할의 학습 인덱스에서만 적합한 표준화 변환을 거친다.
따라서 이후의 모든 계수는 해당 특징이 학습 구간에서 1표준편차 변할 때의 효과로
읽는다.

\paragraph{(F1) 중기 모멘텀}
조정종가 또는 대체 종가 $P_t$로
\begin{equation}
x^{mom}_{i,t}=\log\!\left(\frac{P_t}{P_{t-20}}\right)
\label{eq:momentum}
\end{equation}
와 같이 계산한다. 20거래일은 약 1개월에 해당하는 사전 지정 창으로, 일별 가격
잡음을 완화하면서도 다음 날 행동과 연결되는 중기 추세를 표현한다. 자료에서
성능을 최적화해 고른 값이 아니라 해석 가능성을 위한 설계값이다. 양의 계수는
추세추종을, 음의 계수는 역추세를 뜻하며, 부호 기대는 유형별 모멘텀 거래
성향에 관한 선행 연구~\cite{grinblatt2000investment,choe1999foreign}에서
도출된다.

\paragraph{(F2) 군집행동}
유형 $i$를 제외한 두 유형 $j,k$의 전일 행동에 대한 횡단면 평균
\begin{equation}
x^{herd}_{i,t}=\frac{u_{j,t-1}+u_{k,t-1}}{2},\qquad j,k\neq i
\label{eq:herd}
\end{equation}
으로 정의한다. 시간평균이 아니라 같은 시점의 다른 두 유형에 대한 평균임에
유의한다. 20일 이동평균을 쓰면 즉각적 반응이 아니라 한 달 누적 수급 압력을
측정하게 되므로, 본 연구는 전일 반응에 초점을 맞춘다. 양의 계수는 다른 유형의
전일 방향을 따르는 행동을, 음의 계수는 반대 방향 행동을 뜻한다. 부호 기대는
기관 군집행동 문헌~\cite{lakonishok1992impact}에서 도출되나, 해당 문헌이 기관
\emph{내부}의 군집을 측정하는 반면 식~\eqref{eq:herd}은 \emph{유형 간} 관계를
측정한다는 점에서 대상이 다르다(\S\ref{sec:construct}).

\paragraph{(F3) 평균단가 대비 손실구간}
유형별 평균단가 대용치 $\bar P_{i,t}$와 당일 총거래 VWAP $P^{vwap}_{i,t}$로
\begin{equation}
x^{uw}_{i,t}=\max\!\left(
\frac{\bar P_{i,t}-P^{vwap}_{i,t}}{\bar P_{i,t}},\,0\right)
\label{eq:underwater}
\end{equation}
와 같이 정의한다. 0에서 절단하므로 이 변수는 거래가격이 평균단가 대용치보다
낮은 구간에서만 작동하며, 그 자체로 준거점 의존적 설계다. 평균단가는 총매수ㆍ
총매도 수량과 금액으로 유효 보유량을 재귀적으로 갱신해 계산하되, 전일 보유량에
망각계수 $\rho=0.98$을 곱해 오래된 거래의 영향을 줄이고 매도는 남은 물량의
평균단가를 바꾸지 않는다. $\rho=0.98$의 반감기는 약 34거래일이고 20거래일 후
약 66.8\%가 남으므로 최근 1--2개월의 거래기억을 강조하며, 이 값 역시 자료에서
최적화하지 않았다. 양의 계수는 평균단가 아래에서 다음 날 매수 강도가 커지는
방향이며, 부호 기대는 처분효과
문헌~\cite{shefrin1985disposition,odean1998investors}에서 도출된다.

이 대용치는 실제 계좌 보유량이 아니라 집계 거래로부터 구성한 근사치다. 같은
유형 안의 서로 다른 계좌가 수행한 매수와 매도가 넷팅되므로 ``대표 투자자의
평균단가''라는 개념 자체가 근사에 불과하며, 따라서 $x^{uw}$의 계수는
처분효과를 직접 검정한 결과로 해석할 수 없다.

\subsection{시장 컨텍스트}
\label{sec:context}

컨텍스트는 두 변수로 구성한다. (C1) KOSPI200 1일 수익률 $r_t$, (C2) 원/달러
환율 수준
\begin{equation}
z_t=\frac{\log e_t-\mu_{t-1,252}}{\sigma_{t-1,252}}
\label{eq:fxcontext}
\end{equation}
이다. 환율의 현재값은 사용하지만 평균과 표준편차는 $t-1$까지의 252거래일만
사용한다. 두 컨텍스트 역시 각 분할의 학습 인덱스에서만 다시 표준화하므로,
계수는 해당 컨텍스트가 학습 구간에서 1표준편차 변할 때의 효과로 읽는다.

컨텍스트와 보상 특징은 서로 독립이 아니다. 삼성전자가 KOSPI200에서 차지하는
비중 때문에 종목 모멘텀과 시장수익률이 겹치며, 특징--컨텍스트 상관은 $-0.29$에서
$0.22$ 범위에 있다(그림~\ref{fig:featcorr}). 따라서 개별 계수는 다른 항이
통제된 조건부 연관이며 단독으로 인과적 크기로 읽어서는 안 된다.

\subsection{근시안적 보상, 정책, 그리고 Lasso로의 환원}
\label{sec:model}

\paragraph{잠재 행동점수}
유형 $i$의 시점별 행동점수를
\begin{equation}
q_{i,t}=(\beta_i+B_iC_t)^{\top}x_{i,t}+\alpha_i^{\top}C_t
\label{eq:score}
\end{equation}
로 둔다. $\beta_i\in\mathbb{R}^{3}$은 컨텍스트가 평균일 때의 기본 보상 가중치,
$B_i\in\mathbb{R}^{3\times 2}$는 특징--컨텍스트 상호작용,
$\alpha_i\in\mathbb{R}^{2}$는 컨텍스트의 직접 주효과다. $\alpha$항은 종목
특징이 평균값일 때도 시장 국면이 행동 수준을 이동시키는 수준효과이고, $B$항은
국면에 따라 각 특징에 대한 민감도가 달라지는 기울기 효과다. 상호작용항을
포함하면서 하위차수 항 $C_t$를 누락하는 것은 표준적 사양 오류이므로,
$\alpha$의 포함은 자료로 정당화할 대상이 아니라 사양상 요구되는 것이다.

\paragraph{근시안적 오목 보상}
연속 행동 $a$에 대한 단기 보상을
\begin{equation}
R_i(x,C,a)=q_{i,t}\,a-\frac{\kappa}{2}a^{2}
\label{eq:reward}
\end{equation}
로 둔다. 이차항은 큰 순매수와 순매도에 대칭적 비용을 부과하고 보상을 행동에
대해 오목하게 만들어 내부해를 가능하게 한다. 이 항이 없으면 $[-1,1]$에서의
최대화 해가 $q$의 부호에 따라 거의 항상 $\pm1$이 되어 행동 \emph{강도}를
재현할 수 없다. 식~\eqref{eq:reward}을 $a$에 대해 최대화하면 닫힌형 정책
\begin{equation}
a^{*}_{i,t+1}=\mathrm{clip}\!\left(\frac{q_{i,t}}{\kappa},-1,1\right)
\label{eq:policy}
\end{equation}
을 얻는다. 곡률 $\kappa$는 식별을 위해 1로 고정한다. $\kappa=2$이면 같은 $q$에서
행동 크기가 $1/2$이 되지만 재학습 시 $q$의 계수를 2배로 바꾸면 같은 행동이
재현되므로, $\kappa=1$은 보상과 계수의 단위를 고정하는 정규화이지 곡률이
경험적으로 1이라는 주장이 아니다.

\paragraph{추정}
유형별 학습목적은 학습집합 $\mathcal{D}_i$에서
\begin{equation}
\min_{\theta_i}\ \frac{1}{|\mathcal{D}_i|}
\sum_{t\in\mathcal{D}_i}\bigl(a^{*}_{i,t+1}-a_{i,t+1}\bigr)^{2}
+\lambda_i\lVert\theta_i\rVert_1
\label{eq:loss}
\end{equation}
이며, Adam~\cite{kingma2015adam}으로 유형별 모형을 독립 학습한다.

\begin{proposition}[Lasso 환원]
\label{prop:lasso}
$z_{i,t}=[\,x_{i,t};\ \mathrm{vec}(x_{i,t}C_t^{\top});\ C_t\,]\in\mathbb{R}^{11}$,
$\theta_i=[\,\beta_i;\ \mathrm{vec}(B_i);\ \alpha_i\,]$로 두면
식~\eqref{eq:score}는 $q_{i,t}=\theta_i^{\top}z_{i,t}$이다. 학습 표본 전체에서
$|q_{i,t}|\le 1$이면 식~\eqref{eq:policy}의 clip이 작동하지 않아 $a^{*}=q$이고,
식~\eqref{eq:loss}은
\begin{equation}
\min_{\theta_i}\ \frac{1}{|\mathcal{D}_i|}
\sum_{t}\bigl(\theta_i^{\top}z_{i,t}-a_{i,t+1}\bigr)^{2}
+\lambda_i\lVert\theta_i\rVert_1
\label{eq:lasso}
\end{equation}
이 된다. 이는 다음 날 순매수 강도를 확장 상태특징 $z$에 회귀하는 $L_1$
회귀(Lasso)와 정확히 같다.
\end{proposition}

즉 $\kappa=1$의 근시안적 오목 보상 아래에서 ``선호를 복원한다''와 ``수급을
회귀한다''는 수치적으로 구별되지 않는다. 그림~\ref{fig:saturation}은
명제~\ref{prop:lasso}의 전제가 본 표본에서 실제로 성립함을 보인다. 세 유형 모두
잠재 행동점수의 분포가 행동 경계 $\pm1$로부터 크게 떨어져 있으며($|q|$의
최대값은 외국인 0.842, 개인 0.587, 기관 0.135), 45개 분할 전부에서 포화 비율이
정확히 0이다.

\paragraph{이 환원을 숨기지 않는 이유}
이 일치는 방법의 약점이 아니라 근시안 가정의 논리적 귀결이다. 에이전트가 한
걸음만 내다본다고 가정하면 오늘의 행동이 내일 이후에 영향을 주지 않으므로 내다볼
미래가 없고, 따라서 무엇을 원하는가(보상)와 무엇을 하는가(정책)가 하나로
겹친다. 최대엔트로피 IRL~\cite{ziebart2008maximum}의 확률적 정책
$\pi(a\mid x,C)\propto\exp(R/\tau)$를 식~\eqref{eq:reward}에 적용하면 $\pi$는
평균 $q$, 분산 $\tau$의 절단정규분포가 되어 로그가능도 최대화의 점추정치가
식~\eqref{eq:lasso}과 일치하므로, 이 환원은 특정 구현의 우연이 아니라
근시안ㆍ이차보상 설정 전반의 성질이다.

역강화학습 정식화의 구별되는 가치는 확장 경로에 있다. 상태 동역학과 다기간
가치를 도입하는 순간 보상은 정책과 갈라지고, 회귀로는 도달할 수 없는 구조적
복원으로 이어진다. 본 논문은 그 경로의 첫 단계이며, 성립 조건을 명시한 위에
해석과 검증을 더한다.

\subsection{검증 프로토콜}
\label{sec:protocol}

복원된 보상을 계수 부호와 크기의 보고로 끝내지 않고 다음 네 단계를 적용한다.

\begin{description}
\item[(V1) CPCV 표본외 안정성] 10개 시간 폴드 중 2개를 시험 폴드로 선택하는
조합형 정화 교차검증~\cite{lopez2018advances}을 사용한다. 조합 수는
$\binom{10}{2}=45$이며, 학습과 시험 사이에 1일 purge와 5일 embargo를 적용하고
모든 표준화 변환은 각 분할의 학습 인덱스에만 적합한다. 45개 분할에서 지배 부호가
나타난 분할의 비율(부호 일관성)을 보고하되, 분할들이 관측을 공유하므로 부호
일관성은 독립표본 $p$-값이 아니라 안정성 지표로만 해석한다.

\item[(V2) 달력월 블록 부트스트랩] 부호 일관성이 검정이 아니라는 한계를
보완하기 위해, 표본을 달력월 단위 블록으로 나누고 월 블록을 복원추출해 1{,}000회
재추정하여 계수의 95\% 신뢰구간을 구한다. 월 블록을 쓰는 이유는 일별 수급의
자기상관과 월중 패턴을 블록 내부에 보존하기 위해서다.

\item[(V3) Expanding walk-forward] CPCV는 시험 폴드가 학습 폴드보다 앞설 수
있으므로, 시간 방향을 엄격히 지키는 검증을 추가한다. 최초 400개 관측으로 학습해
5일 purge 후 다음 60거래일을 예측하고, 학습 구간을 60거래일씩 확장하며
반복한다(총 10개 창).

\item[(V4) 정규화 방식 교체] 결과가 $L_1$ 벌점의 산물이 아님을 확인하기 위해
동일 프로토콜을 $L_2$ 벌점과 무벌점으로 재실행하고 계수 부호를 비교한다.
\end{description}

\paragraph{평가지표}
행동 강도 오차는 MAE와 RMSE로, 방향 재현은
$\mathrm{sign}(a^{*})=\mathrm{sign}(a)$의 비율인 방향 정확도로 평가한다. 실제
행동과 예측 행동의 피어슨 상관, $|a^{*}|=1$인 포화 비율을 함께 보고한다. 유형 간
행동 표준편차가 다르므로 MAE와 RMSE의 절대값은 유형 간 비교에 쓰지 않고 각 유형
안에서 해석하며, 표준편차로 정규화한 오차를 병기한다. 참조 기준선으로 전일
자기행동을 그대로 예측하는 지속성 규칙을 함께 보고한다(\S\ref{sec:notclaim}).

\paragraph{구성 타당성 점검}
복원된 부호를 독립적으로 확립된 행동재무 규칙성에서 도출한 사전 기대와
대조한다(\S\ref{sec:construct}). 이 대조는 사후 해석이 아니라 사전에 정한 부호
기대와의 비교다.

% =====================================================================
\section{실험}
\label{sec:experiments}

\subsection{데이터와 설정}
\label{sec:data}

원자료는 삼성전자(005930) 일별 자료이며 가격, 거래대금, 유형별 총매수ㆍ총매도
금액과 수량, 원/달러 환율, KOSPI200 수익률을 포함한다. 252일 환율 이력, 20일
모멘텀, 다음 날 행동을 모두 계산할 수 있는 관측 중 최신 973행을 분석 표본으로
사용하였다. 상태일은 2022년 1월 6일--2025년 12월 29일이며, 마지막 상태는
2025년 12월 30일 행동을 목표로 한다.

단일 종목은 표본의 한계가 아니라 교차종목 교란 없이 유형 간 선호 이질성만
분리하기 위한 통제된 시연 설정이다. 종목 고정효과, 유동성 차이, 지수 편입 효과가
개입하지 않으므로 유형 간 계수 차이를 설계가 아닌 학습의 결과로 귀속할 수 있다.

\begin{table}[t]
  \caption{데이터 및 추정 설정}
  \label{tab:setup}
  \small
  \begin{tabular}{@{}>{\raggedright\arraybackslash}p{0.30\columnwidth}
                  >{\raggedright\arraybackslash}p{0.62\columnwidth}@{}}
    \toprule
    항목 & 설정 \\
    \midrule
    표본 & 973거래일, 2022-01-06--2025-12-29 \\
    유형 & 외국인, 기관, 개인 (대칭 추정) \\
    보상 특징 & 모멘텀(20일), 군집(전일), 손실구간($\rho=0.98$) \\
    컨텍스트 & KOSPI200 1일 수익률, USD/KRW 252일 수준 $z$ \\
    유형별 파라미터 & 11개 ($\beta$ 3, $B$ 6, $\alpha$ 2) \\
    공통 설정 & seed 42, CPU, $\kappa=1$, 학습구간 전용 표준화 \\
    외국인 & 75 epochs, batch 256, lr 0.001, $\lambda=0$ \\
    기관 & 10 epochs, batch 2048, lr 0.0005, $\lambda=0.01$ \\
    개인 & 20 epochs, batch 512, lr 0.001, $\lambda=0.0003$ \\
    검증 & CPCV 10 folds, test 2, purge 1, embargo 5, 45 splits \\
    안정성 & 월 블록 부트스트랩 1{,}000회, walk-forward 10창, 정규화 교체 \\
    \bottomrule
  \end{tabular}
\end{table}

행동 평균은 외국인 $-0.0063$, 기관 $-0.0100$, 개인 $0.0035$로 세 유형 모두 0에
가까우나 표준편차는 개인(0.334)이 가장 크고 기관(0.125)이 가장 작다. 1차
자기상관은 외국인 0.401, 기관 0.149, 개인 0.345다.

동일 거래일 유형 간 행동은 강한 음의 관계를 갖는다(그림~\ref{fig:crossinv}).
외국인--개인 상관이 $-0.772$, 기관--개인이 $-0.499$이며, 이는 한 유형의 매수가
다른 유형의 매도로 나타날 수밖에 없는 시장 청산 제약을 반영한다. 이 사실은
\S\ref{sec:beta}과 \S\ref{sec:alpha}의 해석 범위를 제한하는 근거로 다시
사용된다.

\subsection{표본외 행동 재현}
\label{sec:oos}

\begin{figure}[t]
  \centering
  \includegraphics[width=\columnwidth]{../figures_v2/fig2_actual_vs_prediction.pdf}
  \caption{투자자별 실제 행동과 CPCV 표본외 예측 행동. 각 날짜의 9개 CPCV
  시험예측을 먼저 평균한 뒤 실제값과 예측값에 20거래일 이동평균을 적용했다.
  괄호 안 상관은 이동평균 계열 간 값이며 표~\ref{tab:performance}의 일별
  표본외 상관과 직접 비교하지 않는다.}
  \Description{외국인, 기관, 개인의 실제 순매수 행동과 모형 예측을 시간축 위에
  겹쳐 그린 세 개의 선그래프. 기관 패널에서 예측선이 0 근방에 평탄하게 놓여 있다.}
  \label{fig:prediction}
\end{figure}

그림~\ref{fig:prediction}은 세 유형의 실제 순매수 행동과 표본외 예측 행동을
나란히 보인다. 세 가지가 한눈에 읽힌다. 첫째, 외국인과 개인에서는 예측 계열이
실제 계열의 국면 전환을 따라간다(이동평균 계열 간 상관 0.69, 0.76). 둘째, 예측
진폭이 실제보다 뚜렷하게 작으며, 이는 $L_1$ 정규화에 의한 계수 축소의 결과다.
셋째, 기관의 예측 계열은 사실상 0선에 붙어 있다(0.02). \S\ref{sec:null}에서
보고하는 null이 계수 표가 아니라 예측 시계열에서 먼저 드러난다.

\begin{table}[t]
  \caption{CPCV 표본외 성능(45개 분할 평균 $\pm$ 표준편차)과 지속성 기준선}
  \label{tab:performance}
  \small
  \centering
  \begin{tabular}{@{}lccc@{}}
    \toprule
    지표 & 외국인 & 기관 & 개인 \\
    \midrule
    \multicolumn{4}{@{}l}{\emph{주모형}}\\
    방향 정확도   & $0.636\pm0.039$ & $0.542\pm0.038$ & $0.608\pm0.044$ \\
    상관계수      & $0.278\pm0.110$ & $0.068\pm0.064$ & $0.245\pm0.118$ \\
    MAE           & $0.195\pm0.019$ & $0.094\pm0.008$ & $0.264\pm0.029$ \\
    RMSE          & $0.246\pm0.021$ & $0.125\pm0.012$ & $0.318\pm0.030$ \\
    MAE / 행동 SD & $0.759$ & $0.756$ & $0.791$ \\
    포화 비율     & $0.000$ & $0.000$ & $0.000$ \\
    \addlinespace
    \multicolumn{4}{@{}l}{\emph{지속성 기준선} ($a_{t+1}\leftarrow a_t$)}\\
    방향 정확도   & $\mathbf{0.642}$ & $\mathbf{0.547}$ & $\mathbf{0.616}$ \\
    상관계수      & $\mathbf{0.359}$ & $\mathbf{0.136}$ & $\mathbf{0.311}$ \\
    MAE           & $0.223$ & $0.123$ & $0.299$ \\
    \addlinespace
    행동의 1차 자기상관 & $0.401$ & $0.149$ & $0.345$ \\
    \bottomrule
  \end{tabular}
\end{table}

\begin{figure}[t]
  \centering
  \includegraphics[width=\columnwidth]{../figures_v2/fig5_oos_performance.pdf}
  \caption{표본외 성능의 분할 간 분포와 자명한 기준선. 실선은 주모형 평균,
  파선은 지속성 기준선이다. 세 유형 모두에서 파선이 실선보다 위에 있다.}
  \Description{방향 정확도와 상관계수에 대해 세 투자자 유형의 45개 분할 값을
  바이올린 및 점으로 그린 두 개의 패널.}
  \label{fig:oos}
\end{figure}

표~\ref{tab:performance}과 그림~\ref{fig:oos}에서 평균만 보면 가려지는 두 가지가
드러난다. 첫째, 분할 간 산포가 크다. 특히 상관계수는 외국인이 $-0.05$에서
$0.45$까지 걸쳐 있어 단일 분할의 수치로 유형을 비교하는 것은 의미가 없다. 둘째,
자명한 지속성 기준선이 방향 정확도와 상관계수에서 주모형을 앞선다. $a_t$는 $t$일
종가 이후 공개되므로 이는 정보상 이용 가능한 규칙이다. 주모형이 앞서는 것은
MAE뿐인데, 이는 정규화로 예측 진폭이 축소되어 크기 오차가 줄어든 결과이므로
정보량 우위로 읽을 수 없다.

이 격차의 원인은 명확하다. 행동의 1차 자기상관이 $0.401/0.149/0.345$인데 본
연구의 보상 특징 집합은 자기 지연행동을 포함하지 않는다. 특징을 (F1)--(F3)로
제한한 것은 세 유형이 독립적으로 확립된 행동재무 규칙성에 대응하는 특징만
공유하도록 설계했기 때문이며, 예측 성능을 극대화하려는 설계가 아니다. 본 연구는
이 결과를 은폐하지 않고 표~\ref{tab:performance}에 함께 싣는다.
표~\ref{tab:performance}과 그림~\ref{fig:oos}는 모형이 잡음만 학습한 것이 아님을
보이는 용도로 사용되며, 예측 성능의 우위를 주장하는 근거로 사용되지 않는다.

모든 유형ㆍ모든 분할에서 포화 비율이 정확히 0이므로 명제~\ref{prop:lasso}의 전제
$|q|\le1$이 표본 전체에서 성립한다. 즉 본 추정은 실제로 식~\eqref{eq:lasso}의
Lasso와 동일하며, 성능이 $\pm1$ 극단 예측에 의존하지도 않는다.

\subsection{복원된 보상: 모멘텀 축의 거울상}
\label{sec:beta}

\begin{figure*}[t]
  \centering
  \includegraphics[width=\textwidth]{../figures_v2/fig1_reward_weights_bar.pdf}
  \caption{투자자별로 복원된 기본 보상가중치 $\beta$. 파란색은 양, 붉은색은 음의
  평균 계수이며, 오차막대는 45개 분할의 1 표준편차, 괄호 안은 부호 일관성이다.}
  \Description{외국인, 기관, 개인 각각에 대해 모멘텀, 군집, 손실구간 계수를 그린
  세 개의 막대그래프. 외국인 모멘텀은 양, 개인 모멘텀은 음이며 기관의 세 막대는
  거의 보이지 않는다.}
  \label{fig:betabar}
\end{figure*}

\begin{figure*}[t]
  \centering
  \includegraphics[width=\textwidth]{../figures_v2/fig3_beta_alpha_forest.pdf}
  \caption{보상가중치와 컨텍스트 주효과의 분할 간 안정성. 회색 점은 개별 CPCV
  분할, 굵은 오차선은 45개 분할의 표준편차, 옅은 가로선은 달력월 블록 부트스트랩
  95\% 신뢰구간이다. 속이 빈 표식은 신뢰구간이 0을 포함하는 항이다.}
  \Description{패널 A는 기본 보상가중치, 패널 B는 컨텍스트 주효과에 대한 forest
  plot으로, 각 항의 분할별 추정치와 신뢰구간을 보인다.}
  \label{fig:forest}
\end{figure*}

그림~\ref{fig:betabar}는 세 유형에 대해 복원된 기본 보상가중치를 나란히 보인다.
동일한 특징 집합과 동일한 모형 형태를 부여하고 파라미터만 분리 학습했으므로, 세
패널의 높이 차이는 전적으로 데이터에서 학습된 것이다.
그림~\ref{fig:forest}는 같은 $\beta$와 컨텍스트 주효과 $\alpha$를 분할 수준까지
펼쳐 안정성을 보인다.

\paragraph{모멘텀 축의 거울상}
외국인의 모멘텀 가중치는 $+0.0496$으로 45개 분할 전부에서 양수였고, 개인은
$-0.0304$로 45개 분할 전부에서 음수였다. 그림~\ref{fig:betabar}에서 두 막대가
0선을 사이에 두고 반대 방향을 향하며, 그림~\ref{fig:forest}에서 두 부트스트랩
신뢰구간은 서로 겹치지 않고 각각 0을 제외한다. 동일한 특징 집합을 유형별로 독립
추정했는데도 두 유형이 모멘텀에서 정반대로 갈린다. 외국인은 추세추종, 개인은
역추세로 나타나며, 이는 선행 연구~\cite{grinblatt2000investment,choe1999foreign}의
부호 기대와 일치한다.

\paragraph{군집}
군집 계수는 세 유형 모두 45개 분할 전부에서 음수였고, 외국인($-0.0402$)과
개인($-0.0322$)은 부트스트랩 신뢰구간도 0을 제외한다. 그러나 이 결과는 행동적
해석과 시장 청산 제약을 분리하지 못한다. 그림~\ref{fig:crossinv}가 보이듯 동일
거래일 유형 간 행동 상관이 $-0.77$에 이르므로, 음의 군집 계수 중 상당 부분이
회계적 제약에서 발생할 수 있다. 식~\eqref{eq:herd}가 1거래일 지연값을 쓰므로
동시점 제약이 직접 적용되지는 않으나, 행동의 자기상관을 통해 제약이 전이될 수
있다. 따라서 본 연구는 이를 ``관측된 집계 수급이 유형 간 음의 교차 연관을
보인다''는 서술적 결과로 제한하고, 각 유형이 독립적으로 반대매매 성향을 갖는다는
행동적 주장으로 확장하지 않는다.

\paragraph{손실구간}
개인의 손실구간 가중치는 $+0.0293$으로 45개 분할 전부에서 양수였다. 평균단가
대용치 아래에서 다음 날 매수 강도가 커지는 방향이며, 처분효과
문헌~\cite{shefrin1985disposition,odean1998investors}의 기준가격 의존성과
정합한다. 외국인도 $+0.0094$(95.6\%)로 같은 부호이나 크기가 개인의 $1/3$
수준이다.

다만 그림~\ref{fig:forest}에서 보듯 손실구간은 세 유형 모두 부트스트랩 신뢰구간이
0을 포함한다. 개인의 구간은 $[-0.0122,+0.0665]$로, 45개 분할에서 부호가 100\%
일관됨에도 불구하고 월 블록 재표집에서는 부호가 뒤집히는 경우가 5.5\% 발생한다.
이는 CPCV 분할들이 관측을 공유하기 때문에 부호 일관성이 불확실성을 과소평가한다는
\S\ref{sec:protocol}의 경고가 실제로 작동한 사례다. 따라서 본 연구는 처분효과와
정합하는 부호가 관측되었다고 보고하되, 이를 통계적으로 확립된 결과로 주장하지
않는다.

\subsection{컨텍스트 반응: 거울상의 재현}
\label{sec:alpha}

모멘텀 축에서 관찰된 거울상이 컨텍스트 주효과에서도 그대로
재현된다(그림~\ref{fig:forest}의 패널 B). 외국인은 시장 상승일에 다음 날 순매수
방향이 강해지고($+0.0361$) 원화 약세 국면에서는 순매도 방향으로
이동하는($-0.0261$) 반면, 개인은 두 컨텍스트 모두에서 정확히
반대였다($-0.0246$, $+0.0298$). 네 계수 모두 45개 분할 전부에서 부호가
유지되었고 부트스트랩 신뢰구간도 모두 0을 제외한다. 기관의 두 계수는 부호
일관성이 95\% 이상이나 절대 크기가 외국인ㆍ개인의 $1/8$ 수준이고 신뢰구간이 0을
포함한다.

이 거울상 역시 군집 계수와 같은 해석상 유보가 필요하다.
그림~\ref{fig:crossinv}의 $-0.772$라는 상관 하에서 두 유형의 컨텍스트 반응이
반대 부호로 나타나는 것은 부분적으로 시장 청산 제약을 반영할 수 있다. 본 논문은
이를 두 유형이 시장ㆍ환율 국면에서 반대편 수급을 형성한다는 서술적 사실로
제시한다.

\paragraph{상호작용}
12개 상호작용 계수 중 45개 분할 전부에서 부호가 유지된 것은 외국인의
모멘텀 $\times$ KOSPI200($-0.0139$), 개인의 모멘텀 $\times$ 환율($+0.0144$),
개인의 군집 $\times$ KOSPI200($-0.0080$), 개인의 손실구간 $\times$
환율($-0.0186$) 네 개이며, 이 중 부트스트랩을 통과하는 것은 앞의 하나와 마지막 하나다. 12개 계수를
동시에 검토했으므로 본 연구는 상호작용을 탐색적 결과로만 제시한다.

\subsection{안정성 프로토콜의 결과}
\label{sec:stability}

\paragraph{(V2) 달력월 블록 부트스트랩}
그림~\ref{fig:forest}의 옅은 가로선이 그 결과다. 부트스트랩은 CPCV 부호
일관성보다 엄격한 기준이며, 결과는 셋으로 갈린다. 첫째, 모멘텀 거울상과 컨텍스트
거울상은 부트스트랩에서도 유지된다. 외국인ㆍ개인의 모멘텀, 군집, 두 컨텍스트
주효과 등 8개 항이 모두 0을 제외한다. 둘째, 손실구간 가중치는 통과하지 못한다.
셋째, 기관은 어떤 항도 0을 제외하지 못하며, 그림~\ref{fig:forest}에서 기관의
여섯 표식이 모두 속이 비어 있다.

\begin{figure*}[t]
  \centering
  \includegraphics[width=\textwidth]{../figures_v2/fig6_walk_forward.pdf}
  \caption{Expanding walk-forward 10개 창의 계수 경로와 방향 정확도. 성능은
  창마다 흔들리지만 계수 부호는 유지된다.}
  \Description{확장 창 번호를 가로축으로 하여 세 유형의 모멘텀 계수, 두 컨텍스트
  주효과, 방향 정확도를 그린 네 개의 선그래프.}
  \label{fig:walkforward}
\end{figure*}

\paragraph{(V3) Expanding walk-forward}
그림~\ref{fig:walkforward}가 그 결과다. 두 가지가 동시에 읽힌다. 성능은 창마다
크게 흔들린다. 외국인 방향 정확도는 0.43에서 0.85 사이를 오가며 평균
$0.610\pm0.116$, 상관계수는 $0.131\pm0.162$로 CPCV의 0.278에서 절반 이하로
떨어진다. 개인도 $0.245\rightarrow0.129$다. CPCV는 시험 폴드를 학습 폴드보다
앞세울 수 있는 반면 walk-forward는 그렇지 않으므로, 본 연구는 이 저하를 시간
구조에서 오는 일반화 한계로 그대로 보고한다.

반면 계수 경로는 흔들리지 않는다. 외국인의 모멘텀 계수는 10개 창 전부에서 양수,
개인은 전부에서 음수이며 두 경로가 한 번도 교차하지 않는다. 두 컨텍스트
주효과에서도 같은 거울상이 창 전체에 걸쳐 유지된다. 기관의 세 경로는 모두 0선에
붙어 있다. 즉 성능은 시간에 따라 흔들리지만 복원된 선호의 방향은 흔들리지
않는다.

\begin{figure*}[t]
  \centering
  \includegraphics[width=\textwidth]{../figures_v2/fig7_institution_null.pdf}
  \caption{기관의 수렴된 null. 좌: 복원된 선호의 크기(로그 축). 우: 정규화 방식
  교체에 대한 민감도이며 세 패널은 각각 독립 $y$축을 쓴다.}
  \Description{왼쪽은 유형별 계수 절대값을 로그 축 막대로, 오른쪽은 L1, L2,
  무벌점 하의 계수를 유형별로 그린 세 개의 선그래프.}
  \label{fig:null}
\end{figure*}

\paragraph{(V4) 정규화 방식 교체}
그림~\ref{fig:null}이 그 결과다. 외국인과 개인의 계수는 벌점 방식에 사실상
불변이다(소수점 넷째 자리까지 일치). 따라서 모멘텀 거울상은 $L_1$ 벌점의 산물이
아니다. 반면 기관의 계수만 벌점에 따라 이동하며(군집
$-0.0076\rightarrow-0.0124$, 손실구간 $-0.0000\rightarrow-0.0018$), 그 $y$축
범위는 다른 두 유형의 $1/20$ 수준이다. 계수가 벌점 선택에 좌우된다는 것은
데이터가 그 값을 결정하지 못한다는 뜻이다.

\subsection{기관: 수렴된 null}
\label{sec:null}

기관 수급에 대해 본 연구는 일별 빈도에서 안정적으로 식별되는 선호가 없다는
결과를 얻었고, 이를 은폐하지 않고 발견으로 보고한다. 근거는 다섯이며 모두 그림에
나타난다.

\begin{enumerate}
\item \textbf{예측의 퇴화.} 그림~\ref{fig:prediction} 중간 패널에서 기관의 예측
계열은 실제 계열이 $\pm0.1$ 사이를 오가는 동안 0선에 붙어 있다.
\item \textbf{크기.} 그림~\ref{fig:betabar}에서 기관 패널의 세 막대는 다른 두
유형과 같은 $y$축 위에서 거의 보이지 않는다. 군집을 제외하면 모멘텀 $-0.0006$,
손실구간 $-0.0008$로 $1/40$ 수준이다.
\item \textbf{불확실성.} 그림~\ref{fig:forest}에서 기관의 여섯 표식은 모두 속이
비어 있다.
\item \textbf{벌점 의존성.} 그림~\ref{fig:null} 우측 중앙 패널에서 기관만 벌점에
따라 계수가 이동한다.
\item \textbf{시간 불안정.} 그림~\ref{fig:walkforward}에서 기관의 경로는 0선
근방에 붙어 있으며, 모멘텀 부호 일관성은 10개 창 중 40\%, 손실구간은 0\%다.
\end{enumerate}

\begin{figure*}[t]
  \centering
  \includegraphics[width=\textwidth]{../figures_v2/fig4_training_convergence.pdf}
  \caption{45개 CPCV 분할에서의 학습 수렴. 옅은 회색선은 개별 분할, 굵은 선과
  띠는 평균 $\pm$ 1 표준편차다. 주석은 도달한 해가 동일 목적함수의 폐쇄형
  최적해 대비 남긴 상대 격차다.}
  \Description{에포크를 가로축으로 하여 세 유형의 훈련 MSE 곡선을 그린 세 개의
  패널. 기관 패널의 감소폭이 가장 작다.}
  \label{fig:convergence}
\end{figure*}

이 null은 최적화 실패가 아니다. 그림~\ref{fig:convergence}는 45개 분할의 학습
곡선을 유형별로 보인다. 기관의 훈련 MSE는 학습 구간 전체에서 1.5\%만 감소하며
곡선이 시각적으로 평탄해 보이지 않는다. 그러나 이는 학습 부족이 아니라 애초에
줄일 수 있는 오차가 거의 없기 때문이다. 명제~\ref{prop:lasso}에 따라
식~\eqref{eq:loss}은 볼록 문제이므로 전역 최적해가 유일하며, 도달한 해와 동일
목적함수의 폐쇄형 최적해를 비교하면 기관의 목적함수 상대 격차는 평균 0.17\%,
최대 0.48\%에 불과하다(외국인 0.21\%, 개인 0.92\%). 즉 세 유형 모두 도달 가능한
최적해에 근접했으며, 기관에서만 그 최적해의 계수가 0 근방이다. 데이터가 기관에
대해 결정하는 선호가 없는 것이지, 최적화가 그것을 찾지 못한 것이 아니다.

경제적으로 이 null은 자연스럽게 읽힌다. ``기관''은 연기금, 투신, 보험, 사모,
금융투자를 합산한 범주이므로 서로 다른 목적함수를 가진 주체들의 수급이 상쇄된다.
유일하게 일관된 음의 군집 계수는 이들이 공통적으로 다른 유형의 반대편에 서는
유동성 공급자 역할을 한다는 해석과 정합하나, 부트스트랩 신뢰구간이 경계에
걸치므로 확정하지 않는다.

\subsection{구성 타당성 점검}
\label{sec:construct}

\begin{table}[t]
  \caption{사전 부호 기대와 복원 결과}
  \label{tab:construct}
  \small
  \begin{tabular}{@{}>{\raggedright\arraybackslash}p{0.26\columnwidth}
                  >{\raggedright\arraybackslash}p{0.15\columnwidth}
                  >{\raggedright\arraybackslash}p{0.22\columnwidth}
                  >{\raggedright\arraybackslash}p{0.24\columnwidth}@{}}
    \toprule
    규칙성 & 사전 기대 & 복원 결과 & 판정 \\
    \midrule
    외국인의 추세거래~\cite{grinblatt2000investment,choe1999foreign}
      & $\beta^{mom}>0$ & $+0.0496$ (100\%, CI 0 제외) & 부합 \\
    개인의 역추세 거래~\cite{grinblatt2000investment,kaniel2008individual}
      & $\beta^{mom}<0$ & $-0.0304$ (100\%, CI 0 제외) & 부합 \\
    처분효과~\cite{shefrin1985disposition,odean1998investors}
      & $\beta^{uw}>0$ & $+0.0293$ (100\%, CI 0 포함) & 부분 부합 \\
    기관의 군집행동~\cite{lakonishok1992impact}
      & $\beta^{herd}>0$ & $-0.0044$ (100\%, CI 0 포함) & 불일치 \\
    \bottomrule
  \end{tabular}
\end{table}

표~\ref{tab:construct}이 대조 결과다. 세 유형에 동일한 특징과 동일한 모형 형태를
부여하고 파라미터만 분리 학습했음에도, 복원된 부호가 각각 독립적으로 확립된
규칙성과 부합하고 두 유형이 모멘텀에서 정반대로 갈린다. 이는 복원된 보상이 임의의
적합 결과가 아니라는 구성 타당성의 증거다.

불일치 항목은 측정 대상이 다르기 때문이다. Lakonishok
외~\cite{lakonishok1992impact}는 기관 \emph{내부} 주체들이 같은 방향으로 몰리는
정도를 측정하는 반면, 식~\eqref{eq:herd}은 기관이 \emph{다른 두 유형}의 전일
수급과 같은 방향으로 움직이는지를 측정한다. 따라서 본 결과는 해당 문헌의 반증이
아니며, 유형 내부 군집을 검정하려면 기관 세부 주체별 자료가 필요하다.

\subsection{본 연구가 주장하지 않는 것}
\label{sec:notclaim}

프로토콜의 마지막 단계로, 본 결과가 지지하지 않는 주장을 명시한다.

\begin{enumerate}
\item \textbf{수익률 예측가능성.} 본 연구의 지표는 수급의 방향과 강도에 대한
재현 오차이며 수익률, 샤프지수, 거래비용을 포함하지 않는다.
표~\ref{tab:performance}의 수치로부터 어떤 매매 규칙도 도출되지 않는다.

\item \textbf{수급 예측 성능의 우위.} 표~\ref{tab:performance}과
그림~\ref{fig:oos}에서 보고했듯 지속성 기준선이 방향 정확도와 상관계수에서
주모형을 앞선다. 본 연구의 목적은 예측 성능이 아니라 해석 가능한 공통 척도에서의
선호 복원이며, 두 목적은 다른 특징 집합을 요구한다. 예비 점검에서 지연
자기행동을 네 번째 특징으로 추가하면 모형이 기준선을 넘어서는 반면(방향 정확도
$0.646/0.569/0.629$), 개인의 $\alpha^{KOSPI}$ 부호가 반전되고 손실구간 계수가
소멸한다. 즉 본 논문이 보고한 거울상 중 일부는 수급 자기상관과 분리되지 않는다.

\item \textbf{인과효과.} 복원된 가중치는 표준화된 관측 특징 사이의 조건부
연관이다. 그림~\ref{fig:featcorr}이 보이듯 컨텍스트와 보상 특징 사이에 $-0.29$에서
$0.22$ 범위의 상관이 존재하므로 개별 계수의 크기는 다른 항의 통제 조건에
의존한다.

\item \textbf{유형 간 음의 연관의 행동적 해석.} 군집 계수와 거울상 컨텍스트
반응은 시장 청산 제약과 분리되지 않는다(그림~\ref{fig:crossinv}).

\item \textbf{처분효과의 검정.} 개인의 손실구간 부호는 문헌과 정합하나 월 블록
부트스트랩을 통과하지 못했고(그림~\ref{fig:forest}), 평균단가 대용치는 계좌별
실현손익과 보유기간이 아니라 집계 거래에서 구성한 근사치다.

\item \textbf{종목 간 일반화.} 단일 종목 설정은 유형 간 이질성 분리를 위한
통제이며, 다른 종목과 시장으로의 확장은 검증되지 않았다.
\end{enumerate}

% =====================================================================
\appendix

\section{보충 자료}
\label{sec:appendix}

본문에서 참조한 보조 그림과 표를 모은다. 그림~\ref{fig:saturation}은
명제~\ref{prop:lasso}의 전제를, 그림~\ref{fig:featcorr}은 특징과 컨텍스트의
상관을, 그림~\ref{fig:crossinv}은 유형 간 동시점 행동을 보인다.
표~\ref{tab:interactions}은 본문에서 탐색적 결과로만 제시한 상호작용 계수
전체이고, 표~\ref{tab:convergence}은 도달한 해의 최적성 점검이다.

\begin{figure}[htbp]
  \centering
  \includegraphics[width=\columnwidth]{../figures_v2/figA1_no_saturation.pdf}
  \caption{잠재 행동점수의 분포와 행동 경계. $|q|\le1$이므로
  명제~\ref{prop:lasso}의 전제가 표본 전체에서 성립하며, 45개 분할의 포화 비율은
  정확히 0이다.}
  \Description{세 유형의 잠재 행동점수 히스토그램과 $\pm1$ 경계선.}
  \label{fig:saturation}
\end{figure}

\begin{figure}[htbp]
  \centering
  \includegraphics[width=\columnwidth]{../figures_v2/figA2_feature_correlation.pdf}
  \caption{보상 특징과 컨텍스트의 상관행렬(외국인 특징 기준, 973개 관측).}
  \Description{다섯 변수 사이의 상관계수를 색과 숫자로 표시한 히트맵.}
  \label{fig:featcorr}
\end{figure}

\begin{figure}[htbp]
  \centering
  \includegraphics[width=\columnwidth]{../figures_v2/figA3_cross_investor_actions.pdf}
  \caption{동일 거래일 유형 간 행동. 점선은 완전 상쇄선이며, 외국인과 개인의
  행동이 그 근방에 밀집한다.}
  \Description{유형 쌍별 행동 산점도 세 개와 각 쌍의 상관계수.}
  \label{fig:crossinv}
\end{figure}

\begin{table}[htbp]
  \caption{상호작용 계수 $B$(45개 분할)}
  \label{tab:interactions}
  \small
  \centering
  \begin{tabular}{@{}llrrr@{}}
    \toprule
    유형 & 특징 $\times$ 컨텍스트 & 평균 & 표준편차 & 부호 일관성 \\
    \midrule
    외국인 & 모멘텀 $\times$ KOSPI200 & $-0.0139$ & $0.0052$ & 100.0\% \\
    외국인 & 모멘텀 $\times$ 환율     & $-0.0030$ & $0.0057$ & 71.1\% \\
    외국인 & 군집 $\times$ KOSPI200   & $+0.0016$ & $0.0074$ & 66.7\% \\
    외국인 & 군집 $\times$ 환율       & $+0.0037$ & $0.0073$ & 68.9\% \\
    외국인 & 손실구간 $\times$ KOSPI200 & $-0.0062$ & $0.0040$ & 93.3\% \\
    외국인 & 손실구간 $\times$ 환율   & $-0.0075$ & $0.0076$ & 86.7\% \\
    \addlinespace
    기관 & 모멘텀 $\times$ KOSPI200 & $-0.0001$ & $0.0005$ & 66.7\% \\
    기관 & 모멘텀 $\times$ 환율     & $+0.0003$ & $0.0006$ & 80.0\% \\
    기관 & 군집 $\times$ KOSPI200   & $-0.0004$ & $0.0006$ & 86.7\% \\
    기관 & 군집 $\times$ 환율       & $+0.0012$ & $0.0012$ & 93.3\% \\
    기관 & 손실구간 $\times$ KOSPI200 & $+0.0028$ & $0.0012$ & 97.8\% \\
    기관 & 손실구간 $\times$ 환율   & $+0.0007$ & $0.0008$ & 91.1\% \\
    \addlinespace
    개인 & 모멘텀 $\times$ KOSPI200 & $+0.0021$ & $0.0042$ & 80.0\% \\
    개인 & 모멘텀 $\times$ 환율     & $+0.0144$ & $0.0045$ & 100.0\% \\
    개인 & 군집 $\times$ KOSPI200   & $-0.0080$ & $0.0046$ & 100.0\% \\
    개인 & 군집 $\times$ 환율       & $+0.0092$ & $0.0081$ & 82.2\% \\
    개인 & 손실구간 $\times$ KOSPI200 & $-0.0022$ & $0.0046$ & 66.7\% \\
    개인 & 손실구간 $\times$ 환율   & $-0.0186$ & $0.0046$ & 100.0\% \\
    \bottomrule
  \end{tabular}
\end{table}

\begin{table}[htbp]
  \caption{도달한 해의 최적성 점검(45개 분할)}
  \label{tab:convergence}
  \small
  \centering
  \begin{tabular}{@{}lccc@{}}
    \toprule
    유형 & 훈련 MSE & 목적함수 격차(\%) & 코사인 \\
     & 감소율 & 평균 / 최대 & 평균 / 최소 \\
    \midrule
    외국인 & 14.1\% & $0.21/3.10$ & $0.994/0.948$ \\
    기관   & 1.5\%  & $0.17/0.48$ & $0.935/0.832$ \\
    개인   & 10.3\% & $0.92/1.40$ & $0.894/0.833$ \\
    \bottomrule
  \end{tabular}
\end{table}


\bibliographystyle{ACM-Reference-Format}
\bibliography{paper-ko,paper-refs-additional}

\end{document}
